\documentclass[12pt,a4paper,oneside]{article}

\usepackage[spanish]{babel}
\usepackage[left=2.54cm,right=2.54cm,top=2.54cm,bottom=2.54cm]{geometry}
\usepackage{fancyhdr}
\usepackage{xcolor}
\usepackage{graphicx}
\usepackage{amsmath}
\usepackage{amssymb}

\usepackage{enumitem}
\usepackage{setspace}
\usepackage[hidelinks]{hyperref}
\usepackage{apacite}
\usepackage{multirow}
\usepackage{tabularx} % En el preámbulo

\doublespacing
\setlength{\parindent}{1.27cm}

% Definición de colores basados en la paleta corporativa
\definecolor{primary}{HTML}{227C91}
\definecolor{secondary}{HTML}{605952}
\definecolor{accent}{HTML}{ECEBE9}
\definecolor{destructive}{HTML}{EF4444}
\definecolor{background}{HTML}{FAFAFA}
\definecolor{muted}{HTML}{F1F0EE}
\definecolor{mutedforeground}{HTML}{736B5E}
\definecolor{border}{HTML}{E5E1DD}

% Definición de datos institucionales
\newcommand{\Institucion}{Universidad Nacional Agraria de la Selva}
\newcommand{\Facultad}{Facultad de Ingeniería en Informática y Sistemas}
\newcommand{\Curso}{DISEÑO DETALLADO DE SOFTWARE}
\newcommand{\Tema}{Manual de Identidad Visual y Experiencia de Usuario para el Sistema FINGEST}
\newcommand{\Docente}{Acevedo Aliaga, Alverto}
\newcommand{\FechaEntrega}{28 de octubre de 2025}
\newcommand{\AutorUno}{Evaristo Jara, Jose Junior}
\newcommand{\AutorDos}{Tapullima Serna, Meselemias}
\newcommand{\AutorTres}{Carhuapoma Peña, Nilver Leimer}
\newcommand{\AutorCuatro}{Rengifo Fretel, Juan Manuel}
\newcommand{\AutorCinco}{Castro Cano, Brennis Benjamin}
\pagestyle{fancy}
\fancyhf{}
\rhead{FINGEST - Guía de Identidad}
\lhead{\thepage}
\renewcommand{\headrulewidth}{0.5pt}

\begin{document}

% Portada según normas APA
\begin{titlepage}
  \centering
  \vspace*{0cm}
  {\scshape\large \Institucion \par}
  {\scshape\large \Facultad \par}

  % --- Dos logos centrados ---
  \vspace{1cm}
  \includegraphics[width=3cm]{unas-logo.png}
  \hspace{3.5cm}
  \includegraphics[width=4cm]{fiis-logo.png}
  \vspace{1cm}

  {\scshape\large \Curso \par}
  \vspace{1cm}
  {\bfseries\Large \Tema \par}
  \vspace{1cm}
  \vspace{0.5cm}
  {\large \textbf{Docente:}\par}
  {\large \Docente \par}
  
  \vspace{1cm}
  
  {\large \textbf{Integrantes:} \par}
  \vspace{0.5cm}
  \begin{tabular}{p{6cm} p{6cm}}
    \AutorUno & \AutorDos \\[0.4cm]
    \AutorTres & \AutorCuatro \\[0.4cm]
    \multicolumn{2}{c}{\AutorCinco}
  \end{tabular}
  
  \vfill
  {\large \FechaEntrega \par}
\end{titlepage}

\begin{abstract}
\noindent Este documento presenta el manual de identidad visual y experiencia de usuario para FINGEST, un sistema de gestión financiera digital. Se establecen directrices basadas en principios teóricos de percepción visual (Gestalt), usabilidad (Nielsen, 1994) y accesibilidad (WCAG 2.1) para garantizar coherencia en todos los productos digitales de la organización. El manual define la paleta cromática, sistema tipográfico, componentes de interfaz, y patrones de interacción que conforman la identidad visual del sistema. Los resultados incluyen especificaciones técnicas para implementación en plataformas web y móviles, asegurando una experiencia consistente y accesible para todos los usuarios.

\textbf{Palabras clave:} Identidad visual, experiencia de usuario, diseño de interfaces, accesibilidad web, sistemas de diseño
\end{abstract}

\newpage
\tableofcontents
\newpage

\section{Introducción}

La identidad visual corporativa constituye un elemento fundamental en la construcción de marca y la experiencia de usuario en sistemas digitales (Wheeler, 2017). Este manual establece las directrices de diseño para FINGEST, un sistema de gestión financiera que requiere coherencia visual y funcional en todas sus implementaciones digitales. El documento presenta de manera sistemática los elementos constitutivos de la identidad corporativa: logotipo, sistema cromático, tipografía, principios de usabilidad y aplicaciones prácticas, garantizando la coherencia visual necesaria para fortalecer el reconocimiento y posicionamiento de la marca.

Según Nielsen (1994), la consistencia en interfaces de usuario reduce la carga cognitiva y mejora significativamente la eficiencia operativa. Los principios de Gestalt (Wertheimer, 1923) proporcionan un marco teórico para la organización visual efectiva de elementos en pantalla. La integración de estos fundamentos teóricos con las directrices de accesibilidad WCAG 2.1 (W3C, 2018) asegura que el sistema sea utilizable por el mayor número posible de usuarios, independientemente de sus capacidades.

Este manual está dirigido a diseñadores UI/UX responsables de crear interfaces consistentes, desarrolladores front-end encargados de implementar componentes según especificaciones técnicas, gerentes de producto que toman decisiones estratégicas alineadas con la identidad visual, equipos de calidad que validan el cumplimiento de estándares, y profesionales de marketing y comunicación que producen materiales promocionales. El alcance de aplicación abarca todos los productos digitales desarrollados bajo la marca FINGEST: aplicaciones web y móviles desarrolladas en React, sitios institucionales, interfaces administrativas internas, extensiones de terceros, y materiales de comunicación digital como correos electrónicos y presentaciones.

La correcta aplicación de estas directrices contribuye a consolidar una imagen corporativa sólida, profesional y distintiva que refleja los valores de confianza, innovación y accesibilidad que caracterizan a FINGEST en el mercado de servicios financieros digitales.

\section{Marco Teórico}

Esta sección presenta los fundamentos conceptuales y científicos que sustentan las decisiones de diseño implementadas en FINGEST. La combinación de teorías perceptuales, principios de usabilidad y estándares de accesibilidad proporciona una base sólida para la construcción de interfaces efectivas y centradas en el usuario.

\subsection{Principios de Gestalt en Diseño de Interfaces}

Los principios de Gestalt, desarrollados por psicólogos alemanes en la década de 1920 (Wertheimer, 1923; Koffka, 1935), explican cómo el cerebro humano organiza elementos visuales en patrones coherentes. Estos principios son fundamentales para crear interfaces intuitivas y fáciles de procesar cognitivamente.

El principio de proximidad establece que los elementos cercanos entre sí se perciben como relacionados funcionalmente (Todorovic, 2008). En FINGEST, este principio guía la agrupación de campos de formulario relacionados, la organización de métricas en el panel principal, y la disposición de filtros en las interfaces de búsqueda. Por ejemplo, los campos de correo electrónico y contraseña en el formulario de acceso se posicionan próximos para indicar su relación funcional.

El principio de similitud indica que elementos con características visuales similares se perciben como parte de un grupo (Palmer, 1992). La implementación consistente de estilos visuales para botones primarios, el uso uniforme de iconografía en la navegación, y el tratamiento homogéneo de tarjetas de contenido refuerzan este principio. La continuidad y el cierre completan el marco perceptual, facilitando la comprensión de relaciones jerárquicas y la completación mental de formas parcialmente visibles.

\subsection{Heurísticas de Usabilidad de Nielsen}

Jakob Nielsen (1994) estableció diez principios heurísticos que han servido como referencia fundamental para la evaluación y diseño de interfaces usables. Estas heurísticas proporcionan criterios específicos para crear sistemas que los usuarios puedan operar eficientemente con mínima frustración.

La visibilidad del estado del sistema requiere que los usuarios siempre sepan qué está ocurriendo mediante retroalimentación apropiada y oportuna. FINGEST implementa notificaciones toast para confirmar acciones exitosas, indicadores de carga durante procesos asíncronos, y badges visuales para comunicar estados de elementos. La correspondencia entre el sistema y el mundo real asegura que el lenguaje, conceptos e iconos utilizados sean familiares para los usuarios, empleando términos financieros reconocibles en lugar de jerga técnica.

El control y libertad del usuario se manifiesta en la capacidad de deshacer acciones, cancelar procesos en curso, y navegar libremente entre secciones sin restricciones artificiales. La consistencia y estándares se logran mediante el uso sistemático de componentes de interfaz reutilizables, implementación uniforme de patrones de interacción, y adherencia a convenciones establecidas de diseño web. La prevención de errores mediante validaciones proactivas, confirmaciones para acciones destructivas, y diseño de formularios que minimizan posibilidades de error complementan estas heurísticas fundamentales.

\subsection{Estándares de Accesibilidad WCAG}

Las Pautas de Accesibilidad para el Contenido Web (WCAG) 2.1, publicadas por el World Wide Web Consortium (W3C, 2018), establecen criterios técnicos para hacer contenido web accesible a personas con diversas discapacidades. Estos estándares no solo cumplen requisitos legales, sino que mejoran la experiencia para todos los usuarios.

El principio de perceptibilidad requiere que la información sea presentada de manera que todos los usuarios puedan percibirla, independientemente de sus capacidades sensoriales. FINGEST cumple con ratios de contraste mínimos de 4.5:1 (nivel AA) para texto normal y 7:1 (nivel AAA) para elementos críticos, proporciona texto alternativo para todas las imágenes funcionales, y utiliza múltiples indicadores visuales más allá del color para comunicar información importante.

La operabilidad asegura que todos los elementos interactivos sean accesibles mediante teclado, con indicadores de foco visibles y orden de tabulación lógico. La comprensibilidad se logra mediante lenguaje claro, comportamiento predecible de componentes, y mensajes de error específicos y constructivos. Finalmente, la robustez garantiza compatibilidad con tecnologías asistivas mediante el uso apropiado de HTML semántico y atributos ARIA cuando sea necesario.

\section{Elementos de Identidad Visual}

Esta sección detalla los componentes visuales fundamentales que conforman la identidad de marca de FINGEST. Los elementos aquí descritos deben aplicarse consistentemente en todas las implementaciones digitales para asegurar reconocimiento inmediato y coherencia estética. Cada componente ha sido diseñado considerando tanto principios estéticos como funcionales de accesibilidad y usabilidad.

\subsection{Sistema Cromático}

La paleta de colores constituye uno de los elementos más reconocibles de una identidad visual (Labrecque \& Milne, 2012). En sistemas de gestión financiera, la selección cromática debe transmitir profesionalismo, confianza y estabilidad, mientras cumple con estrictos requisitos de accesibilidad.

FINGEST utiliza una paleta cromática fundamentada en el modelo HSL (Hue, Saturation, Lightness) que facilita ajustes sistemáticos de luminosidad para garantizar contraste adecuado. El color primario (\#227C91, HSL: 191° 62\% 35\%, RGB: 34, 124, 145) transmite confianza y profesionalismo asociado con tonalidades azul-verde utilizadas frecuentemente en el sector financiero. Este color se utiliza en botones de acción principal, enlaces de navegación, y elementos interactivos primarios. El análisis de contraste muestra un ratio de 7.8:1 sobre fondo blanco, superando el nivel AAA de WCAG 2.1.

El color secundario (\#605952, HSL: 30° 8\% 35\%, RGB: 96, 89, 82) proporciona neutralidad cromática para texto secundario y elementos de soporte, con un ratio de contraste de 7.2:1 sobre fondo blanco. El color destructivo (\#EF4444, HSL: 0° 84\% 60\%, RGB: 239, 68, 68) se reserva exclusivamente para alertas, errores y acciones irreversibles, cumpliendo con el ratio mínimo AA de 4.7:1. Los colores complementarios (muted, accent, border) completan el sistema con roles específicos en estados de interfaz, fondos de elementos deshabilitados, y delimitación de áreas de contenido.

\subsection{Tipografía}
La tipografía en interfaces digitales afecta directamente la legibilidad, jerarquía visual y percepción de profesionalismo (Shaikh, 2007). La selección de familias tipográficas debe priorizar la claridad en múltiples resoluciones de pantalla y diversos contextos de visualización.

FINGEST adopta Inter como familia tipográfica principal, una fuente sans-serif diseñada específicamente para interfaces digitales con optimización en renderizado subpixel (Rasmus Andersson, 2020). Inter ofrece excepcional legibilidad en tamaños pequeños, amplio soporte de caracteres Unicode, múltiples pesos (400, 500, 600, 700) para jerarquía visual, y características OpenType para mejora de legibilidad. La tipografía se carga desde Google Fonts mediante la directiva \texttt{@import} en la hoja de estilos global, asegurando disponibilidad sin requerir instalación local.

\subsubsection{Sistema de Jerarquía Tipográfica}

El sistema tipográfico establece una jerarquía clara mediante la variación de tamaño, peso y altura de línea. La Tabla \ref{tab:tipografia} presenta las especificaciones completas del sistema tipográfico de FINGEST.

\begin{table}[h]
\centering
\caption{Especificaciones del Sistema Tipográfico FINGEST}
\label{tab:tipografia}
\begin{tabular}{|l|c|c|c|p{4cm}|}
\hline
\textbf{Elemento} & \textbf{Tamaño} & \textbf{Peso} & \textbf{Line Height} & \textbf{Uso Principal} \\
\hline
H1 & 32px & 700 & 1.2 (38.4px) & Títulos principales de página \\
\hline
H2 & 24px & 700 & 1.3 (31.2px) & Delimitación de subsecciones \\
\hline
H3 & 20px & 600 & 1.4 (28px) & Agrupación de contenido \\
\hline
P (Párrafo) & 16px & 400 & 1.5 (24px) & Texto de cuerpo principal \\
\hline
Small & 14px & 400 & 1.5 (21px) & Información complementaria \\
\hline
Caption & 12px & 400 & 1.4 (16.8px) & Anotaciones y metadatos \\
\hline
\end{tabular}
\end{table}

Los títulos principales (H1) utilizan 32px con peso 700 y line-height 1.2, optimizados para captar atención en secciones principales. Los títulos secundarios (H2) emplean 24px con peso 700 y line-height 1.3 para delimitar subsecciones. Los subtítulos (H3) se configuran a 20px con peso 600 y line-height 1.4. El texto de párrafo estándar utiliza 16px con peso 400 y line-height 1.5, dimensiones que investigaciones demuestran como óptimas para lectura prolongada en pantalla (Bernard et al., 2003). Texto pequeño (14px) y texto de ayuda (12px) completan la jerarquía para información complementaria y anotaciones mínimas.

\subsubsection{Gestión de Elementos HTML}

\paragraph{Títulos Principales (H1)}
Los elementos H1 representan el nivel más alto de jerarquía visual en la aplicación. Se implementan con las siguientes especificaciones:

\begin{verbatim}
h1 {
  font-size: 32px;
  font-weight: 700;
  line-height: 1.2;
  font-family: 'Inter', sans-serif;
  color: #1a73e8;
  margin-bottom: 24px;
}
\end{verbatim}

\textbf{Criterios de aplicación:}
\begin{itemize}
    \item Uso exclusivo para el título principal de cada vista o página
    \item Máximo un elemento H1 por página para mantener jerarquía semántica
    \item Color principal de marca (\texttt{\#1a73e8}) para reforzar identidad visual
    \item Margen inferior de 24px para separación visual clara
\end{itemize}

\textbf{Ejemplos de uso:} ``Bienvenido a FINGEST'', ``Panel de Control'', ``Gestión de Transacciones''.

\paragraph{Títulos Secundarios (H2)}
Los elementos H2 estructuran las secciones principales dentro de una vista:

\begin{verbatim}
h2 {
  font-size: 24px;
  font-weight: 700;
  line-height: 1.3;
  font-family: 'Inter', sans-serif;
  color: #1f2937;
  margin-top: 32px;
  margin-bottom: 16px;
}
\end{verbatim}

\textbf{Criterios de aplicación:}
\begin{itemize}
    \item Delimitan secciones principales dentro de una página
    \item Pueden aparecer múltiples H2 por página para organizar contenido
    \item Color de texto primario (\texttt{\#1f2937}) para mantener legibilidad
    \item Espaciado superior de 32px y inferior de 16px para jerarquía visual
\end{itemize}

\textbf{Ejemplos de uso:} ``Resumen Financiero'', ``Transacciones Recientes'', ``Configuración de Cuenta''.

\paragraph{Párrafos (P)}
Los elementos de párrafo constituyen el contenido de lectura principal:

\begin{verbatim}
p {
  font-size: 16px;
  font-weight: 400;
  line-height: 1.5;
  font-family: 'Inter', sans-serif;
  color: #4b5563;
  margin-bottom: 16px;
  max-width: 600px;
}
\end{verbatim}

\textbf{Criterios de aplicación:}
\begin{itemize}
    \item Contenido de lectura principal y descripciones extensas
    \item Color de texto secundario (\texttt{\#4b5563}) para reducir fatiga visual
    \item Margen inferior de 16px entre párrafos consecutivos
    \item Ancho máximo recomendado de 65-75 caracteres (aprox. 600px) para legibilidad óptima según estudios de usabilidad
\end{itemize}

\paragraph{Contenedores (DIV)}
Los elementos DIV funcionan como contenedores estructurales sin estilo tipográfico intrínseco:

\begin{verbatim}
div {
  font-size: 16px;
  font-weight: 400;
  line-height: 1.5;
  font-family: 'Inter', sans-serif;
}
\end{verbatim}

Los contenedores DIV heredan propiedades tipográficas del elemento padre y se modifican mediante clases utility específicas para variaciones contextuales. Las principales clases aplicables son: \texttt{.text-sm} (14px) para texto reducido, \texttt{.text-xs} (12px) para anotaciones, \texttt{.font-semibold} (peso 600) para énfasis moderado, y \texttt{.font-bold} (peso 700) para énfasis fuerte.

\subsubsection{Sistema de Clases Utility (text-size)}

El sistema implementa clases utility predefinidas para facilitar la gestión tipográfica consistente. La Tabla \ref{tab:text-classes} detalla las clases disponibles y sus aplicaciones.

\begin{table}[h]
\centering
\caption{Clases Utility de Tamaño Tipográfico}
\label{tab:text-classes}
\begin{tabular}{|l|c|c|p{5cm}|}
\hline
\textbf{Clase CSS} & \textbf{Tamaño} & \textbf{Line Height} & \textbf{Aplicación Práctica} \\
\hline
.text-xs & 12px & 1.4 & Timestamps, tooltips, badges \\
\hline
.text-sm & 14px & 1.5 & Labels, botones secundarios, tablas \\
\hline
.text-base & 16px & 1.5 & Texto predeterminado, inputs, cards \\
\hline
.text-lg & 18px & 1.5 & Texto destacado, valores numéricos \\
\hline
.text-xl & 20px & 1.4 & Equivalente a H3 en componentes \\
\hline
.text-2xl & 24px & 1.3 & Equivalente a H2 en componentes \\
\hline
.text-3xl & 32px & 1.2 & Equivalente a H1 en componentes \\
\hline
\end{tabular}
\end{table}

\subsubsection{Consideraciones de Implementación}

\paragraph{Responsividad}
El sistema tipográfico implementa una estrategia de diseño adaptativo mediante breakpoints específicos, garantizando legibilidad óptima en cada contexto de visualización. La Tabla \ref{tab:responsividad} presenta las especificaciones tipográficas para cada categoría de dispositivo.

\begin{table}[h]
\centering
\caption{Adaptación Tipográfica por Dispositivo}
\label{tab:responsividad}
\small
\begin{tabular}{|l|c|c|c|c|}
\hline
\textbf{Elemento} & \textbf{Desktop} & \textbf{Tablet} & \textbf{Móvil L} & \textbf{Móvil S} \\
 & \textbf{≥1024px} & \textbf{768-1023px} & \textbf{480-767px} & \textbf{<480px} \\
\hline
H1 & 32px / 700 & 28px / 700 & 26px / 700 & 24px / 700 \\
\hline
H2 & 24px / 700 & 22px / 700 & 20px / 700 & 18px / 700 \\
\hline
H3 & 20px / 600 & 18px / 600 & 18px / 600 & 16px / 600 \\
\hline
P (Párrafo) & 16px / 400 & 16px / 400 & 15px / 400 & 14px / 400 \\
\hline
Small & 14px / 400 & 14px / 400 & 13px / 400 & 12px / 400 \\
\hline
Caption & 12px / 400 & 12px / 400 & 11px / 400 & 11px / 400 \\
\hline
\end{tabular}
\end{table}

\textbf{Estrategia de implementación por dispositivo:}

\begin{itemize}
    \item \textbf{Desktop (≥1024px):} Tamaños tipográficos completos según diseño base. Aprovecha el espacio amplio para jerarquías visuales pronunciadas. Line-height estándar (1.2-1.5) para lectura cómoda en monitores de alta resolución. Márgenes laterales generosos (padding horizontal de 48-64px) para evitar líneas excesivamente largas que dificulten el seguimiento visual. Ancho máximo de contenido principal limitado a 1200px para mantener líneas de texto óptimas (65-75 caracteres).
    
    \item \textbf{Tablet (768-1023px):} Reducción moderada del 10-12\% en títulos principales (H1, H2) para optimizar espacio vertical. Mantenimiento de tamaños de párrafo (16px) para lectura sostenida en sesiones extendidas. Ajuste de márgenes laterales a 32-48px. Orientación portrait prioriza verticalidad, optimizando espacios entre secciones (margin-bottom incrementado 20\%). Interacciones táctiles requieren áreas de toque mínimas de 44x44px según Apple Human Interface Guidelines y Material Design en elementos interactivos con texto. Line-height mantenido en 1.5 para párrafos.
    
    \item \textbf{Móvil Landscape (480-767px):} Reducción del 15-20\% en títulos para maximizar contenido visible sin sacrificar jerarquía. Texto de párrafo reducido a 15px para balance entre legibilidad y aprovechamiento de espacio horizontal limitado. Line-height incrementado a 1.6 para compensar tamaños reducidos y mejorar escaneo visual en pantallas pequeñas. Márgenes laterales comprimidos a 20-24px para maximizar ancho útil. Priorización de scroll vertical sobre layouts complejos. Eliminación de sidebars y contenido secundario que compita con el flujo principal.
    
    \item \textbf{Móvil Portrait (<480px):} Reducción agresiva del 20-25\% en títulos principales manteniendo mínimos de legibilidad. Texto de párrafo mínimo de 14px según estudios de usabilidad móvil (Budiu \& Nielsen, 2015) que demuestran umbral crítico de legibilidad. Line-height aumentado a 1.6-1.7 para mejorar escaneo visual y reducir fatiga en lecturas extendidas. Márgenes laterales mínimos de 16px (1rem) para maximizar espacio de contenido. Eliminación de jerarquías visuales complejas en favor de estructuras lineales verticales. Botones y elementos interactivos con padding generoso (12-16px vertical, mínimo 44px altura total) para facilitar interacción táctil precisa. Incremento de espaciado entre elementos interactivos consecutivos (mínimo 8px gap).
\end{itemize}

\textbf{Media queries implementadas:}

\begin{verbatim}
/* Desktop First Approach */
@media (max-width: 1023px) { 
    /* Tablet: Ajustes moderados */
    h1 { font-size: 28px; }
    h2 { font-size: 22px; }
    .container { padding: 0 32px; }
}

@media (max-width: 767px) { 
    /* Mobile Landscape: Optimización horizontal */
    h1 { font-size: 26px; }
    h2 { font-size: 20px; }
    p { font-size: 15px; line-height: 1.6; }
    .container { padding: 0 20px; }
}

@media (max-width: 479px) { 
    /* Mobile Portrait: Máxima compresión legible */
    h1 { font-size: 24px; }
    h2 { font-size: 18px; }
    h3 { font-size: 16px; }
    p { font-size: 14px; line-height: 1.7; }
    .container { padding: 0 16px; }
    button, .interactive-element { 
        min-height: 44px; 
        padding: 12px 16px; 
    }
}
\end{verbatim}

La estrategia desktop-first permite definir el diseño completo como base y aplicar reducciones progresivas mediante media queries, facilitando el mantenimiento del código y asegurando que dispositivos de alta resolución reciban la experiencia tipográfica óptima sin restricciones artificiales.

\paragraph{Accesibilidad}
El sistema tipográfico cumple rigurosamente con las pautas WCAG 2.1 nivel AA, garantizando accesibilidad universal:

\begin{itemize}
    \item \textbf{Contraste de color:} Ratio de contraste mínimo 4.5:1 para texto normal según WCAG 2.1 criterio 1.4.3. Ratio de contraste mínimo 3:1 para texto grande (18pt/24px+ o 14pt/18.5px+ en negrita) según criterio 1.4.3. Implementación verificada mediante herramientas automatizadas (Lighthouse, axe DevTools) y validación manual con analizadores de contraste.
    
    \item \textbf{Espaciado y densidad:} Line-height mínimo de 1.5 para texto de párrafo según WCAG 2.1 criterio 1.4.12, facilitando seguimiento visual de líneas. Espaciado entre párrafos de al menos 2x el tamaño de fuente (32px para párrafos de 16px). Letter-spacing no comprimido por debajo de valores predeterminados del navegador. Word-spacing mínimo de 0.16em para mejorar discriminación de palabras.
    
    \item \textbf{Redimensionamiento de texto:} Soporte para zoom de navegador hasta 200\% sin pérdida de contenido o funcionalidad según criterio 1.4.4. Unidades relativas (rem, em) en lugar de píxeles absolutos para respetar preferencias del usuario. Validación de que el contenido permanece legible y funcional con tamaños de fuente personalizados del sistema operativo.
    
    \item \textbf{Optimización para dislexia:} Fuente Inter seleccionada por sus características de legibilidad mejorada: altura x elevada, contadores abiertos, y diferenciación clara entre caracteres similares (I, l, 1; O, 0). Espaciado entre letras y palabras optimizado para reducir crowding visual. Evitación de texto justificado (text-align: justify) que crea espaciados irregulares problemáticos. Line-height generoso (1.5-1.7) para facilitar tracking visual entre líneas.
    
    \item \textbf{Consideraciones de bajo contraste visual:} Evitación de grises muy claros para texto en fondos blancos (mínimo \#4b5563 para texto secundario). Pruebas de legibilidad con simuladores de deficiencias visuales (deuteranopia, protanopia, tritanopia). Soporte para modo oscuro (dark mode) con ajustes de contraste apropiados manteniendo ratios WCAG.
\end{itemize}

\paragraph{Optimización de Rendimiento}
La implementación tipográfica prioriza la velocidad de carga y experiencia de usuario mediante estrategias específicas:

\begin{itemize}
    \item \textbf{Carga estratégica de fuentes:} Inter se carga desde Google Fonts CDN con la directiva \texttt{font-display: swap}, permitiendo que el texto se muestre inmediatamente con una fuente alternativa del sistema (system-ui, -apple-system, Segoe UI) mientras se descarga Inter. Esto previene el flash invisible de texto (FOIT) y mejora el First Contentful Paint (FCP). Timeout de swap configurado en 100ms para balance óptimo entre performance y experiencia visual.
    
    \item \textbf{Subconjuntos optimizados:} Se cargan únicamente subconjuntos latinos (\texttt{latin, latin-ext}) necesarios para el contexto hispanohablante, reduciendo el peso de descarga de aproximadamente 450KB a 180KB (60\% de reducción). Exclusión de caracteres cirílicos, griegos, asiáticos y otros no utilizados en la aplicación.
    
    \item \textbf{Precarga de recursos críticos:} Los pesos tipográficos críticos (400 Regular y 700 Bold) se precargan en el \texttt{<head>} del documento HTML mediante \texttt{<link rel="preload">} para minimizar el flash de contenido sin estilo (FOUT) y mejorar el Largest Contentful Paint (LCP). Pesos secundarios (500, 600) se cargan de forma asíncrona solo cuando son necesarios.
    
    \item \textbf{Caché y compresión:} Configuración de headers HTTP para caché prolongado (1 año) de archivos de fuentes mediante \texttt{Cache-Control: public, max-age=31536000, immutable}. Compresión WOFF2 (Web Open Font Format 2) con ratio de compresión superior a TTF/OTF (30-50\% más ligero). Soporte de fallback a WOFF para navegadores legacy.
    
    \item \textbf{Fallback stack robusto:} Definición de pila de fuentes alternativas que mantienen métricas similares a Inter para minimizar layout shift durante la carga: \texttt{font-family: 'Inter', -apple-system, BlinkMacSystemFont, 'Segoe UI', Roboto, Oxygen, Ubuntu, sans-serif}. Uso de \texttt{size-adjust}, \texttt{ascent-override}, y \texttt{descent-override} en \texttt{@font-face} para matching preciso de métricas con fuentes del sistema.
    
    \item \textbf{Monitoreo de rendimiento:} Implementación de métricas de Web Vitals para tracking de impacto tipográfico: FCP objetivo <1.8s, LCP objetivo <2.5s, CLS objetivo <0.1. Uso de Resource Timing API para monitorear tiempos de descarga de fuentes y optimizar estrategia de carga según datos reales de usuarios.
\end{itemize}

\subsection{Iconografía}

Los iconos funcionan como elementos visuales que facilitan el reconocimiento rápido de funciones y reducen la dependencia del texto (Huang et al., 2002). Un sistema iconográfico coherente mejora la velocidad de procesamiento visual y la eficiencia operativa del usuario al proporcionar claves visuales inmediatas que aceleran la navegación y comprensión de la interfaz.

FINGEST implementa Lucide React (v0.263.1), una biblioteca de iconos de código abierto caracterizada por su estilo lineal minimalista, stroke uniforme de 2 píxeles, diseño geométrico consistente basado en grid de 24×24px, y catálogo extenso de más de 1000 iconos funcionales. El estilo lineal complementa la estética limpia del sistema visual y permite escalado sin pérdida de claridad en dispositivos de alta densidad de píxeles (Retina, HiDPI). Los iconos se implementan como componentes React SVG nativos, garantizando calidad vectorial en cualquier resolución y permitiendo manipulación dinámica mediante props.

\subsubsection{Especificaciones Técnicas}

\paragraph{Sistema de Tamaños}
Los tamaños de iconos se estandarizan según contexto funcional y jerarquía visual. La Tabla \ref{tab:iconos-tamanios} detalla las especificaciones completas.

\begin{table}[h]
\centering
\caption{Sistema de Tamaños de Iconografía}
\label{tab:iconos-tamanios}
\begin{tabular}{|l|c|c|p{5cm}|}
\hline
\textbf{Contexto} & \textbf{Tamaño} & \textbf{Stroke} & \textbf{Casos de Uso} \\
\hline
Extra Small & 12px & 1.5px & Badges, indicators, metadatos inline \\
\hline
Small & 16px & 2px & Campos de entrada, chips, tags \\
\hline
Base & 18px & 2px & Navegación lateral, menús, listas \\
\hline
Medium & 20px & 2px & Botones secundarios, cards \\
\hline
Large & 24px & 2px & Acciones principales, encabezados, toolbars \\
\hline
Extra Large & 32px & 2px & Hero sections, estados vacíos, onboarding \\
\hline
Jumbo & 48px & 2px & Ilustraciones, páginas de error, marketing \\
\hline
\end{tabular}
\end{table}

\textbf{Implementación en código:}

\begin{verbatim}
import { User, Settings, ChevronRight } from 'lucide-react';

// Tamaño small (16px)
<User size={16} />

// Tamaño base (18px) - default para navegación
<Settings size={18} />

// Tamaño large (24px) - acciones principales
<ChevronRight size={24} strokeWidth={2} />
\end{verbatim}

\paragraph{Sistema Cromático}
El color de los iconos sigue un sistema jerárquico que refuerza la arquitectura de información y facilita el escaneo visual. La Tabla \ref{tab:iconos-colores} presenta la paleta cromática aplicada a iconografía.

\begin{table}[h]
\centering
\caption{Sistema Cromático de Iconografía}
\label{tab:iconos-colores}
\begin{tabular}{|l|c|p{6cm}|}
\hline
\textbf{Estado/Tipo} & \textbf{Color (Hex)} & \textbf{Aplicación} \\
\hline
Principal & \#1a73e8 & Acciones primarias, elementos interactivos activos \\
\hline
Secundario & \#4b5563 & Iconos de navegación, elementos informativos \\
\hline
Terciario & \#9ca3af & Iconos decorativos, elementos deshabilitados \\
\hline
Éxito & \#10b981 & Confirmaciones, estados positivos \\
\hline
Advertencia & \#f59e0b & Alertas, acciones que requieren atención \\
\hline
Error & \#ef4444 & Errores, acciones destructivas \\
\hline
Información & \#3b82f6 & Tooltips, ayuda contextual \\
\hline
Neutral & \#6b7280 & Iconos de utilidad general \\
\hline
Heredado & currentColor & Iconos que adaptan color del contenedor padre \\
\hline
\end{tabular}
\end{table}

Los iconos heredan automáticamente el color del elemento padre mediante la propiedad CSS \texttt{currentColor}, asegurando coherencia cromática sin necesidad de especificación manual. Esta aproximación facilita la gestión de estados hover, focus y active sin duplicación de estilos.

\paragraph{Catálogo de Iconos por Categoría}
El sistema organiza iconos según categorías funcionales para facilitar la selección consistente. La Tabla \ref{tab:iconos-catalogo} presenta el mapeo de categorías a iconos específicos.

\begin{table}[h]
\centering
\caption{Catálogo de Iconos por Categoría Funcional}
\label{tab:iconos-catalogo}
\small
\begin{tabular}{|l|p{8cm}|}
\hline
\textbf{Categoría} & \textbf{Iconos Principales} \\
\hline
Navegación & Home, Menu, ChevronLeft, ChevronRight, ChevronUp, ChevronDown, ArrowLeft, ArrowRight \\
\hline
Acciones & Plus, Minus, Edit, Trash2, Save, Download, Upload, Copy, Share2, ExternalLink \\
\hline
Estados & Check, X, AlertCircle, Info, HelpCircle, Clock, CheckCircle, XCircle, AlertTriangle \\
\hline
Usuario & User, Users, UserPlus, UserMinus, Settings, LogIn, LogOut, Shield, Lock, Unlock \\
\hline
Finanzas & DollarSign, TrendingUp, TrendingDown, PieChart, BarChart3, CreditCard, Wallet, Receipt \\
\hline
Archivo & File, FileText, Folder, FolderOpen, Image, FileSpreadsheet, FilePdf, Paperclip \\
\hline
Comunicación & Mail, MessageSquare, Send, Phone, Video, Bell, BellOff, Inbox \\
\hline
Utilidad & Search, Filter, Calendar, Eye, EyeOff, MoreVertical, MoreHorizontal, RefreshCw \\
\hline
\end{tabular}
\end{table}

\subsubsection{Patrones de Uso}

\paragraph{Iconos en Botones}
Los iconos en botones siguen patrones establecidos para maximizar reconocimiento:

\begin{itemize}
    \item \textbf{Botones con icono + texto:} Icono posicionado a la izquierda del texto con espaciado de 8px. Tamaño de icono: 20px para botones grandes, 18px para botones medianos, 16px para botones pequeños. Color del icono hereda del texto del botón.
    
    \item \textbf{Botones solo icono (icon-only):} Requieren tooltip descriptivo al hacer hover. Tamaño mínimo: 44×44px para cumplir con áreas de toque táctil. Icono centrado vertical y horizontalmente. Padding uniforme de 12px.
    
    \item \textbf{Botones de acción flotante (FAB):} Icono de 24px centrado en botón circular de 56×56px. Sombra elevada para indicar interactividad. Color de icono contrasta con fondo del botón (ratio 4.5:1 mínimo).
\end{itemize}

\textbf{Ejemplo de implementación:}

\begin{verbatim}
// Botón con icono + texto
<button className="btn-primary">
  <Plus size={20} />
  <span>Crear Transacción</span>
</button>

// Botón solo icono con tooltip
<button className="btn-icon" aria-label="Editar">
  <Edit size={20} />
</button>
\end{verbatim}

\paragraph{Iconos en Campos de Entrada}
Los iconos mejoran la comprensión de campos de formulario:

\begin{itemize}
    \item \textbf{Prefijos (leading icons):} Posicionados a la izquierda del input con padding de 12px. Tamaño: 16px. Color: \texttt{\#6b7280} (neutral). Indican el tipo de dato esperado (e.g., Search para búsqueda, Mail para email, Lock para contraseña).
    
    \item \textbf{Sufijos (trailing icons):} Posicionados a la derecha del input. Iconos interactivos (Eye/EyeOff para mostrar/ocultar contraseña) con tamaño 18px. Iconos de estado (Check para validación exitosa, AlertCircle para error) con tamaño 16px.
    
    \item \textbf{Estados de validación:} Icono CheckCircle en verde (\texttt{\#10b981}) para inputs válidos. Icono AlertCircle en rojo (\texttt{\#ef4444}) para inputs con error. Transición suave de 200ms al cambiar estados.
\end{itemize}

\paragraph{Iconos en Navegación}
La navegación lateral utiliza iconos para mejorar escaneo visual:

\begin{itemize}
    \item \textbf{Navegación colapsada:} Solo iconos visibles (24×24px) con tooltip al hover. Espaciado vertical de 8px entre elementos. Indicador de página activa mediante color primario y fondo sutil.
    
    \item \textbf{Navegación expandida:} Icono (18px) + texto con espaciado de 12px. Iconos alineados a la izquierda con margen derecho de 12px. Color heredado del texto del enlace.
    
    \item \textbf{Badges numéricos:} Posicionados en la esquina superior derecha del icono. Diámetro mínimo de 18px. Color de fondo según contexto (rojo para notificaciones urgentes, azul para información).
\end{itemize}

\paragraph{Iconos en Cards y Listas}
Cards y elementos de lista utilizan iconos para categorización visual:

\begin{itemize}
    \item \textbf{Iconos de categoría:} Posicionados en esquina superior izquierda o como prefijo del título. Tamaño: 20-24px. Color según categoría o tipo de contenido.
    
    \item \textbf{Iconos de acción:} Agrupados en esquina superior derecha o al final del card. Tamaño: 18px. Visible solo al hover en desktop, siempre visible en móvil. Espaciado de 8px entre iconos consecutivos.
    
    \item \textbf{Iconos de estado:} Acompañan metadatos como fecha, autor, estado. Tamaño: 14-16px. Color neutral o semántico según contexto.
\end{itemize}

\section{Componentes de Interfaz de Usuario}

Esta sección especifica los elementos interactivos que conforman la interfaz de FINGEST mediante tablas técnicas y especificaciones precisas para facilitar la implementación.

\subsection{Botones}

Los botones comunican jerarquía de acciones mediante diferenciación visual clara (Galitz, 2007). La Tabla \ref{tab:botones-specs} detalla las especificaciones técnicas completas.

\begin{table}[h]
\centering
\caption{Especificaciones Técnicas de Botones}
\label{tab:botones-specs}
\small
\begin{tabular}{|l|c|c|c|c|c|}
\hline
\textbf{Tipo} & \textbf{Fondo} & \textbf{Texto} & \textbf{Borde} & \textbf{Altura} & \textbf{Padding H/V} \\
\hline
Primario & \#227C91 & \#FFFFFF / 500 & none & 44px & 24px / 8px \\
\hline
Secundario & transparent & \#227C91 / 500 & 1px \#E5E1DD & 44px & 24px / 8px \\
\hline
Destructivo & \#EF4444 & \#FFFFFF / 500 & none & 44px & 24px / 8px \\
\hline
Fantasma & transparent & \#227C91 / 500 & none & 40px & 16px / 6px \\
\hline
Pequeño & \#227C91 & \#FFFFFF / 500 & none & 36px & 16px / 6px \\
\hline
\end{tabular}
\end{table}

\begin{table}[h]
\centering
\caption{Estados Interactivos de Botones}
\label{tab:botones-estados}
\small
\begin{tabular}{|l|p{10cm}|}
\hline
\textbf{Estado} & \textbf{Especificación} \\
\hline
Default & Border-radius: 4px, transition: 150ms ease-in-out \\
\hline
Hover & Luminosidad 85\%, cursor: pointer, box-shadow: 0 2px 4px rgba(0,0,0,0.1) \\
\hline
Focus & Outline: 2px solid \#227C91, outline-offset: 2px \\
\hline
Active & Transform: scale(0.98), luminosidad 75\% \\
\hline
Disabled & Opacity: 0.5, cursor: not-allowed, pointer-events: none \\
\hline
Loading & Spinner 16px, texto opacity: 0.7, cursor: wait \\
\hline
\end{tabular}
\end{table}

\textbf{Consideraciones de implementación:}
\begin{itemize}
    \item Área táctil mínima: 44×44px (WCAG 2.1 nivel AA)
    \item Máximo un botón primario por vista
    \item Acciones destructivas requieren confirmación modal
    \item \texttt{aria-label} obligatorio en botones solo icono
\end{itemize}

\subsection{Campos de Entrada}

Puntos críticos de interacción para captura de datos (Jarrett \& Gaffney, 2009). La Tabla \ref{tab:inputs-specs} presenta especificaciones completas.

\begin{table}[h]
\centering
\caption{Especificaciones de Campos de Entrada}
\label{tab:inputs-specs}
\small
\begin{tabular}{|l|c|c|c|c|}
\hline
\textbf{Propiedad} & \textbf{Default} & \textbf{Focus} & \textbf{Error} & \textbf{Success} \\
\hline
Border & 1px \#E5E1DD & 2px \#227C91 & 2px \#EF4444 & 2px \#10B981 \\
\hline
Background & \#FFFFFF & \#FFFFFF & \#FEF2F2 & \#F0FDF4 \\
\hline
Text Color & \#1F2937 & \#1F2937 & \#DC2626 & \#059669 \\
\hline
Placeholder & \#9CA3AF & \#9CA3AF & \#9CA3AF & \#9CA3AF \\
\hline
Border Radius & 4px & 4px & 4px & 4px \\
\hline
Height & 44px & 44px & 44px & 44px \\
\hline
Padding H/V & 12px / 8px & 12px / 8px & 12px / 8px & 12px / 8px \\
\hline
\end{tabular}
\end{table}

\begin{table}[h]
\centering
\caption{Variantes de Campos de Entrada}
\label{tab:inputs-variantes}
\small
\begin{tabular}{|l|p{9cm}|}
\hline
\textbf{Tipo} & \textbf{Especificaciones Adicionales} \\
\hline
Con Icono Prefijo & Icon: 16px, position: left 12px, input padding-left: 40px \\
\hline
Con Icono Sufijo & Icon: 16px, position: right 12px, input padding-right: 40px \\
\hline
Textarea & Min-height: 88px (2 líneas), max-height: 264px, resize: vertical \\
\hline
Select & Chevron icon 18px derecha, padding-right: 40px, appearance: none \\
\hline
Search & Icon: Search 16px izquierda, border-radius: 20px (pill shape) \\
\hline
Number & Arrows ocultos, input type: number, pattern: ``[0-9]*'' \\
\hline
Date & Icon: Calendar 16px derecha, formato: DD/MM/YYYY \\
\hline
\end{tabular}
\end{table}

\textbf{Sistema de validación:}
\begin{itemize}
    \item Mensajes debajo del campo: 12px, color según estado
    \item Campos requeridos: asterisco rojo (*) junto al label
    \item Validación en tiempo real para campos críticos (email, contraseña)
    \item Helper text: 12px \#6B7280, margin-top: 4px
\end{itemize}

\subsection{Tarjetas y Contenedores}

Contenedores visuales que agrupan información relacionada (Krug, 2014). La Tabla \ref{tab:cards-specs} especifica variantes.

\begin{table}[h]
\centering
\caption{Especificaciones de Tarjetas}
\label{tab:cards-specs}
\small
\begin{tabular}{|l|c|c|c|c|}
\hline
\textbf{Tipo} & \textbf{Background} & \textbf{Border} & \textbf{Radius} & \textbf{Shadow} \\
\hline
Base & \#FFFFFF & 1px \#E5E1DD & 8px & 0 1px 2px rgba(0,0,0,0.05) \\
\hline
Elevada & \#FFFFFF & none & 8px & 0 4px 6px rgba(0,0,0,0.1) \\
\hline
Interactiva & \#FFFFFF & 1px \#E5E1DD & 8px & 0 1px 2px rgba(0,0,0,0.05) \\
\hline
Destacada & \#F8FBFF & 1px \#227C91 & 8px & 0 2px 4px rgba(34,124,145,0.1) \\
\hline
\end{tabular}
\end{table}

\begin{table}[h]
\centering
\caption{Anatomía de Tarjetas}
\label{tab:cards-anatomia}
\small
\begin{tabular}{|l|p{9cm}|}
\hline
\textbf{Elemento} & \textbf{Especificaciones} \\
\hline
Padding General & 24px todos los lados \\
\hline
Header & Padding-bottom: 16px, border-bottom: 1px \#E5E1DD (opcional) \\
\hline
Título & Font-size: 20px, weight: 600, color: \#1F2937 \\
\hline
Descripción & Font-size: 14px, weight: 400, color: \#6B7280, margin-top: 4px \\
\hline
Content & Padding-top: 16px (con header), margin-bottom: 16px entre elementos \\
\hline
Footer & Padding-top: 16px, border-top: 1px \#E5E1DD, display: flex, justify: space-between \\
\hline
Actions & Gap: 8px entre botones, align: right \\
\hline
\end{tabular}
\end{table}

\textbf{Estados interactivos (tarjetas clickeables):}
\begin{itemize}
    \item Hover: box-shadow: 0 4px 6px rgba(0,0,0,0.1), transform: translateY(-2px), transition: 200ms
    \item Active: transform: scale(0.98)
    \item Cursor: pointer
\end{itemize}

\subsection{Navegación}

Sistema de navegación lateral y horizontal. La Tabla \ref{tab:nav-specs} detalla especificaciones por breakpoint.

\begin{table}[h]
\centering
\caption{Especificaciones de Navegación Lateral}
\label{tab:nav-specs}
\small
\begin{tabular}{|l|c|c|c|}
\hline
\textbf{Propiedad} & \textbf{Desktop ≥1024px} & \textbf{Tablet 768-1023px} & \textbf{Móvil <768px} \\
\hline
Ancho & 256px (fijo) & 64px (colapsado) & 100\% (drawer) \\
\hline
Posición & Fixed left & Fixed left & Fixed overlay \\
\hline
Background & \#FFFFFF & \#FFFFFF & \#FFFFFF \\
\hline
Border Right & 1px \#E5E1DD & 1px \#E5E1DD & none \\
\hline
Z-index & 1000 & 1000 & 9999 \\
\hline
\end{tabular}
\end{table}

\begin{table}[h]
\centering
\caption{Anatomía de Elementos de Navegación}
\label{tab:nav-items}
\small
\begin{tabular}{|l|p{8cm}|}
\hline
\textbf{Elemento} & \textbf{Especificaciones} \\
\hline
Logo Section & Height: 64px, padding: 16px, border-bottom: 1px \#E5E1DD \\
\hline
Nav Item & Height: 44px, padding: 12px 16px, border-radius: 4px, margin: 4px 8px \\
\hline
Icon & Size: 18px, margin-right: 12px, color: inherit \\
\hline
Label & Font-size: 14px, weight: 500, truncate: ellipsis \\
\hline
Badge & Size: 18px min, border-radius: 9px, font-size: 11px, position: absolute right \\
\hline
Divider & Height: 1px, background: \#E5E1DD, margin: 8px 16px \\
\hline
\end{tabular}
\end{table}

\begin{table}[h]
\centering
\caption{Estados de Elementos de Navegación}
\label{tab:nav-estados}
\small
\begin{tabular}{|l|c|c|c|}
\hline
\textbf{Estado} & \textbf{Background} & \textbf{Text Color} & \textbf{Icon Color} \\
\hline
Default & transparent & \#4B5563 & \#6B7280 \\
\hline
Hover & \#F1F0EE & \#1F2937 & \#227C91 \\
\hline
Active & \#227C91 & \#FFFFFF & \#FFFFFF \\
\hline
Focus & \#F1F0EE & \#1F2937 & \#227C91 \\
\hline
\end{tabular}
\end{table}

\begin{table}[h]
\centering
\caption{Header Horizontal}
\label{tab:header-specs}
\small
\begin{tabular}{|l|p{9cm}|}
\hline
\textbf{Propiedad} & \textbf{Especificación} \\
\hline
Height & 64px fijo \\
\hline
Background & \#FFFFFF \\
\hline
Border Bottom & 1px \#E5E1DD \\
\hline
Padding H & 24px desktop, 16px móvil \\
\hline
Z-index & 999 (debajo de sidebar drawer) \\
\hline
Layout & Flexbox: space-between, align-items: center \\
\hline
Left Section & Logo + Título de página (breadcrumbs opcional) \\
\hline
Right Section & Search + Notifications + User Menu, gap: 16px \\
\hline
\end{tabular}
\end{table}

\textbf{Consideraciones de responsividad:}
\begin{itemize}
    \item Móvil: Menú hamburguesa 44×44px, sidebar drawer con overlay backdrop rgba(0,0,0,0.5)
    \item Tablet: Sidebar colapsado (solo iconos) con tooltip al hover
    \item Desktop: Sidebar expandido persistente, puede colapsarse con toggle
    \item Transiciones: 300ms ease-in-out para colapsar/expandir
\end{itemize}

\section{Patrones de Interacción y Retroalimentación}

Esta sección describe los mecanismos mediante los cuales el sistema comunica estados, procesa acciones del usuario, y proporciona confirmación de operaciones. Los patrones de interacción coherentes reducen la curva de aprendizaje y aumentan la confianza del usuario en el sistema (Norman, 2013).

\subsection{Feedback Visual y Notificaciones}

El feedback visual inmediato es esencial para que los usuarios comprendan el resultado de sus acciones y el estado actual del sistema (Nielsen, 1994). La ausencia de feedback genera incertidumbre y puede llevar a acciones duplicadas o abandono de tareas.

El estado hover en elementos interactivos se indica mediante cambio sutil de color de fondo con transición de 150ms, elevación opcional de sombra en tarjetas clickeables, y cambio de cursor a pointer. El estado click puede incluir efecto de escala a 98\% para simular pulsación física. Los estados de carga se comunican mediante spinners circulares para operaciones de duración indeterminada, barras de progreso para procesos con porcentaje conocido, y skeleton screens para carga inicial de contenido preservando la estructura de layout.

Las notificaciones tipo toast se implementan usando la biblioteca Sonner con posicionamiento en esquina superior derecha, duración de 4 segundos para mensajes informativos y 6 segundos para errores, posibilidad de cierre manual mediante botón X, y clasificación en notificaciones de éxito (fondo verde, icono checkmark), error (fondo destructive, icono de alerta), información (fondo azul, icono info), y advertencia (fondo amarillo, icono de precaución). Los mensajes emplean lenguaje específico y constructivo evitando jerga técnica.

\subsection{Animaciones y Transiciones}

Las animaciones bien ejecutadas mejoran la percepción de fluidez y profesionalismo del sistema, mientras que animaciones excesivas o lentas degradan la experiencia (Laubheimer, 2020). El objetivo es utilizar animación para comunicar relaciones espaciales y temporales, no para decoración arbitraria.

Las micro-interacciones (hover, focus, click) utilizan duraciones de 100-200ms para feedback instantáneo percibido. Las transiciones de página y cambio de vistas emplean máximo 300-500ms para evitar percepción de lentitud. Los tipos de animaciones implementadas incluyen fade in/out para modales, tooltips y notificaciones, slide horizontal o vertical para sidebars, drawers y menús desplegables, scale sutil (0.95-1.05) para botones y tarjetas interactivas, y skeleton loading con efecto shimmer para indicar carga de contenido.

Todas las animaciones deben respetar la preferencia del sistema \texttt{prefers-reduced-motion}, deshabilitándose automáticamente para usuarios que han configurado reducción de movimiento en su sistema operativo, cumpliendo así con la directriz WCAG 2.1 sobre animaciones y movimiento.

\subsection{Accesibilidad e Inclusión}

La accesibilidad no es una característica opcional sino un requisito fundamental que asegura que personas con diversas capacidades puedan utilizar el sistema efectivamente (Henry et al., 2014). El cumplimiento de estándares de accesibilidad además mejora la experiencia para todos los usuarios en diversos contextos de uso.

El contraste de colores cumple con WCAG 2.1 nivel AA mínimo (4.5:1 para texto normal) y nivel AAA cuando es factible (7:1). Los ratios específicos en FINGEST son: primario sobre blanco 7.8:1 (AAA), secundario sobre blanco 7.2:1 (AAA), destructivo sobre blanco 4.7:1 (AA), y muted-foreground sobre muted 4.8:1 (AA). La información nunca se comunica exclusivamente mediante color, utilizando adicionalmente iconos, texto y patrones.

La navegación por teclado asegura que todos los elementos interactivos sean alcanzables mediante tecla Tab, manteniendo orden lógico de tabulación de izquierda a derecha y arriba hacia abajo, mostrando indicadores de focus visibles con anillo de 2px en color primario, y proporcionando atajos de teclado documentados para acciones frecuentes. Los formularios incluyen labels asociados explícitamente mediante atributo \texttt{htmlFor}, mensajes de error vinculados con \texttt{aria-describedby}, instrucciones previas al campo cuando la función no es obvia, y validación en tiempo real que no interfiere con lectores de pantalla.

\section{Adaptación a Dispositivos y Plataformas}

Esta sección establece criterios técnicos para la adaptación de la interfaz mediante tablas de especificación que facilitan la implementación responsiva (Marcotte, 2011).

\subsection{Diseño Responsivo}

FINGEST implementa diseño responsivo con enfoque mobile-first. La Tabla \ref{tab:breakpoints} detalla los breakpoints del sistema.

\begin{table}[h]
\centering
\caption{Sistema de Breakpoints y Características por Dispositivo}
\label{tab:breakpoints}
\small
\begin{tabular}{|l|c|c|p{5cm}|}
\hline
\textbf{Breakpoint} & \textbf{Min Width} & \textbf{Dispositivo} & \textbf{Características Principales} \\
\hline
xs (base) & 0px & Móvil portrait & Diseño vertical, 1 columna, drawer navigation \\
\hline
sm & 640px & Móvil landscape & 1-2 columnas, menú expandible \\
\hline
md & 768px & Tablet portrait & 2 columnas, sidebar colapsado (iconos) \\
\hline
lg & 1024px & Tablet landscape/Laptop & 2-3 columnas, sidebar expandido fijo \\
\hline
xl & 1280px & Desktop & 3-4 columnas, layout completo \\
\hline
2xl & 1536px & Desktop large & 4 columnas, márgenes amplios \\
\hline
\end{tabular}
\end{table}

\subsubsection{Adaptación de Componentes}

La Tabla \ref{tab:componentes-responsive} especifica cómo cada componente se adapta por breakpoint.

\begin{table}[h]
\centering
\caption{Comportamiento Responsivo de Componentes}
\label{tab:componentes-responsive}
\tiny
\begin{tabular}{|l|p{2.5cm}|p{2.5cm}|p{2.5cm}|p{2.5cm}|}
\hline
\textbf{Componente} & \textbf{xs-sm (<768px)} & \textbf{md (768-1023px)} & \textbf{lg (1024-1279px)} & \textbf{xl+ (≥1280px)} \\
\hline
Grid de Cards & 1 columna, gap: 16px & 2 columnas, gap: 20px & 3 columnas, gap: 24px & 4 columnas, gap: 24px \\
\hline
Sidebar & Drawer oculto, overlay & Colapsado (64px), iconos solo & Expandido (256px), fijo & Expandido (280px), fijo \\
\hline
Header & 56px, hamburguesa visible & 64px, hamburguesa visible & 64px, sin hamburguesa & 64px, sin hamburguesa \\
\hline
Formularios & 1 columna, stack vertical & 1 columna, campos anchos & 2 columnas para pares & 2 columnas con sidebar \\
\hline
Tablas & Scroll horizontal o card-stack & Scroll horizontal, columnas ajustadas & Tabla completa, 6-8 columnas & Tabla completa, todas las columnas \\
\hline
Modales & 100\% ancho, padding: 16px & 90\% ancho, max: 600px & 80\% ancho, max: 800px & Max: 1000px, centrado \\
\hline
Botones & Full-width (excepto inline) & Width: auto, min: 120px & Width: auto, padding mayor & Width: auto, padding mayor \\
\hline
Tipografía & H1: 24px, P: 14px & H1: 28px, P: 16px & H1: 32px, P: 16px & H1: 32px, P: 16px \\
\hline
Spacing & Padding: 16px, gap: 12px & Padding: 24px, gap: 16px & Padding: 32px, gap: 24px & Padding: 48px, gap: 24px \\
\hline
\end{tabular}
\end{table}

\subsubsection{Clases Utility Responsivas}

\begin{table}[h]
\centering
\caption{Sistema de Clases Tailwind Responsivas}
\label{tab:tailwind-responsive}
\small
\begin{tabular}{|l|p{9cm}|}
\hline
\textbf{Categoría} & \textbf{Ejemplo de Implementación} \\
\hline
Layout & \texttt{flex flex-col md:flex-row lg:flex-row-reverse} \\
\hline
Grid & \texttt{grid grid-cols-1 md:grid-cols-2 lg:grid-cols-3 xl:grid-cols-4} \\
\hline
Spacing & \texttt{p-4 md:p-6 lg:p-8 xl:p-12} \\
\hline
Typography & \texttt{text-sm md:text-base lg:text-lg} \\
\hline
Display & \texttt{hidden md:block lg:flex} \\
\hline
Width & \texttt{w-full md:w-1/2 lg:w-1/3 xl:w-1/4} \\
\hline
Height & \texttt{h-screen md:h-auto lg:min-h-[600px]} \\
\hline
Position & \texttt{relative md:absolute lg:fixed} \\
\hline
\end{tabular}
\end{table}

\subsubsection{Patrones de Adaptación Específicos}

\paragraph{Navegación Responsiva}

\begin{table}[h]
\centering
\caption{Estados de Navegación por Breakpoint}
\label{tab:nav-responsive}
\small
\begin{tabular}{|l|c|c|c|}
\hline
\textbf{Breakpoint} & \textbf{Sidebar} & \textbf{Header} & \textbf{Interacción} \\
\hline
xs-sm (<768px) & Drawer oculto & Hamburguesa + Logo & Tap hamburguesa → Drawer overlay \\
\hline
md (768-1023px) & Colapsado 64px & Sin hamburguesa & Iconos + Tooltip al hover \\
\hline
lg+ (≥1024px) & Expandido 256px & Sin hamburguesa & Toggle collapse opcional \\
\hline
\end{tabular}
\end{table}

\paragraph{Tablas de Datos Responsivas}

\begin{table}[h]
\centering
\caption{Estrategias para Tablas Responsivas}
\label{tab:tablas-responsive}
\small
\begin{tabular}{|l|p{9cm}|}
\hline
\textbf{Estrategia} & \textbf{Implementación} \\
\hline
Scroll Horizontal & Container: \texttt{overflow-x: auto}, tabla: \texttt{min-width: 600px}, aplicar en xs-sm \\
\hline
Card Stack & Cada fila → Card individual, mostrar solo columnas críticas, aplicar en xs \\
\hline
Columnas Ocultas & \texttt{hidden md:table-cell} en columnas secundarias, priorizar 3-4 columnas en móvil \\
\hline
Accordion Rows & Fila principal + expansión para detalles, icono chevron para expandir \\
\hline
\end{tabular}
\end{table}

\paragraph{Formularios Responsivos}

\begin{table}[h]
\centering
\caption{Layout de Formularios por Breakpoint}
\label{tab:forms-responsive}
\small
\begin{tabular}{|l|p{9cm}|}
\hline
\textbf{Breakpoint} & \textbf{Layout} \\
\hline
xs-sm (<768px) & Stack vertical, labels encima de inputs, botones full-width \\
\hline
md (768-1023px) & 1 columna para campos largos, botones width: auto centrados \\
\hline
lg+ (≥1024px) & 2 columnas para pares (nombre/apellido), grupos relacionados juntos \\
\hline
\end{tabular}
\end{table}

\subsubsection{Testing de Responsividad}

\begin{table}[h]
\centering
\caption{Resoluciones Críticas para Testing}
\label{tab:testing-resolutions}
\small
\begin{tabular}{|l|c|c|p{4cm}|}
\hline
\textbf{Dispositivo} & \textbf{Resolución} & \textbf{Viewport} & \textbf{Prioridad} \\
\hline
iPhone SE & 375×667 & 375px & Alta - Móvil pequeño \\
\hline
iPhone 12/13/14 & 390×844 & 390px & Crítica - Móvil estándar \\
\hline
iPhone Pro Max & 428×926 & 428px & Media - Móvil grande \\
\hline
iPad Mini & 768×1024 & 768px & Alta - Tablet portrait \\
\hline
iPad Pro 11" & 834×1194 & 834px & Media - Tablet grande \\
\hline
Laptop 13" & 1280×800 & 1280px & Crítica - Desktop mínimo \\
\hline
Desktop 1080p & 1920×1080 & 1920px & Alta - Desktop estándar \\
\hline
Desktop 4K & 3840×2160 & 3840px & Media - Desktop alto \\
\hline
\end{tabular}
\end{table}

\textbf{Criterios de validación:}
\begin{itemize}
    \item Legibilidad: Texto mínimo 14px en móvil
    \item Interactividad: Áreas táctiles mínimo 44×44px
    \item Performance: Layout shift (CLS) < 0.1
    \item Funcionalidad: Todas las acciones accesibles sin scroll horizontal
\end{itemize}

\newpage

\subsection{Consideraciones Multiplataforma}

Planificación para futuras implementaciones nativas. La Tabla \ref{tab:plataformas-comparacion} compara características por plataforma.

\begin{table}[h]
\centering
\caption{Especificaciones por Plataforma}
\label{tab:plataformas-comparacion}
\tiny
\begin{tabular}{|l|p{3cm}|p{3cm}|p{3cm}|p{3cm}|}
\hline
\textbf{Aspecto} & \textbf{Web (Actual)} & \textbf{iOS (Futuro)} & \textbf{Android (Futuro)} & \textbf{Desktop (Futuro)} \\
\hline
Design System & Custom + Tailwind & HIG + SwiftUI & Material Design 3 & Electron + Native \\
\hline
Navegación & Sidebar + Header & Tab Bar + Nav Bar & Bottom Nav + Drawer & Sidebar + Menu Bar \\
\hline
Tipografía & Inter 14-32px & SF Pro Dynamic Type & Roboto + System & Inter/System \\
\hline
Colores & Paleta fija & Adaptar a Dark Mode & Material You Theming & Sistema + preferencias \\
\hline
Iconos & Lucide (2px stroke) & SF Symbols (variable) & Material Icons (filled) & Lucide/System \\
\hline
Gestos & Click, scroll básico & Swipe, pinch, 3D Touch & Swipe, long-press & Click, keyboard shortcuts \\
\hline
Notificaciones & Browser push & APNs native & FCM native & Sistema operativo \\
\hline
Almacenamiento & LocalStorage/IndexedDB & CoreData/SwiftData & Room/SQLite & Electron Store \\
\hline
\end{tabular}
\end{table}

\subsubsection{Guías de Diseño por Plataforma}

\paragraph{iOS (Futuro)}

\begin{table}[h]
\centering
\caption{Especificaciones iOS (Human Interface Guidelines)}
\label{tab:ios-specs}
\small
\begin{tabular}{|l|p{9cm}|}
\hline
\textbf{Elemento} & \textbf{Especificación HIG} \\
\hline
Status Bar & 44pt altura, considerar Safe Area, adaptar color a contenido \\
\hline
Navigation Bar & 44pt altura (compact), 96pt (large title), translucent background \\
\hline
Tab Bar & 49pt altura, máximo 5 tabs, iconos 25×25pt, SF Symbols \\
\hline
List Rows & Altura mínima 44pt, swipe actions (leading/trailing) \\
\hline
Buttons & Sistema: SF Pro Text 17pt, prominentes: filled/tinted background \\
\hline
Typography & Dynamic Type: soporte de 11 tamaños, SF Pro (system font) \\
\hline
Spacing & Márgenes: 16-20pt, padding interno: 12-16pt \\
\hline
Haptics & Feedback: impacto, notificación, selección según acción \\
\hline
\end{tabular}
\end{table}

\paragraph{Android (Futuro)}

\begin{table}[h]
\centering
\caption{Especificaciones Android (Material Design 3)}
\label{tab:android-specs}
\small
\begin{tabular}{|l|p{9cm}|}
\hline
\textbf{Elemento} & \textbf{Especificación Material} \\
\hline
Status Bar & 24dp altura, color adaptable a contenido, translúcido \\
\hline
App Bar (Top) & 56dp altura (mobile), 64dp (tablet/desktop), elevación: 0-4dp \\
\hline
Bottom Navigation & 56dp altura, 3-5 items, iconos 24×24dp, Material Icons \\
\hline
FAB & 56×56dp (regular), 40×40dp (mini), elevación: 6dp, margin: 16dp \\
\hline
Cards & Elevación: 1dp (resting), border-radius: 12dp (M3), padding: 16dp \\
\hline
Typography & Roboto/Material 3 Type Scale, soporte de tamaños dinámicos \\
\hline
Colors & Material You: tema dinámico desde wallpaper, tonos adaptativos \\
\hline
Ripple Effect & Touch feedback: bounded ripple, unbounded para iconos \\
\hline
\end{tabular}
\end{table}

\subsubsection{Compatibilidad Cross-Browser (Web)}

\begin{table}[h]
\centering
\caption{Matriz de Compatibilidad de Navegadores}
\label{tab:browser-compatibility}
\small
\begin{tabular}{|l|c|c|c|c|c|}
\hline
\textbf{Característica} & \textbf{Chrome} & \textbf{Firefox} & \textbf{Safari} & \textbf{Edge} & \textbf{Mobile} \\
\hline
CSS Grid & 57+ & 52+ & 10.1+ & 16+ & iOS 10.3+, Android 81+ \\
\hline
Flexbox & 29+ & 28+ & 9+ & 12+ & iOS 9+, Android 4.4+ \\
\hline
Custom Properties & 49+ & 31+ & 9.1+ & 15+ & iOS 9.3+, Android 81+ \\
\hline
ES6 Modules & 61+ & 60+ & 11+ & 16+ & iOS 11+, Android 81+ \\
\hline
Service Workers & 40+ & 44+ & 11.1+ & 17+ & iOS 11.3+, Android 5+ \\
\hline
WebP & 32+ & 65+ & 14+ & 18+ & iOS 14+, Android 4.2+ \\
\hline
Container Queries & 105+ & 110+ & 16+ & 105+ & iOS 16+, Android 105+ \\
\hline
\end{tabular}
\end{table}

\textbf{Estrategias de soporte:}
\begin{itemize}
    \item \textbf{Progressive Enhancement:} Funcionalidad base para todos, características avanzadas para modernos
    \item \textbf{Polyfills:} Cargar condicionalmente para navegadores legacy (Intersection Observer, ResizeObserver)
    \item \textbf{Fallbacks CSS:} \texttt{@supports} para características modernas, alternativas para antiguos
    \item \textbf{Transpilación:} Babel para ES6+ → ES5 en producción para navegadores antiguos
    \item \textbf{Testing:} Automatizado en BrowserStack/Sauce Labs para matriz de navegadores
\end{itemize}


\subsubsection{Optimización por Plataforma}

\begin{table}[h]
\centering
\caption{Optimizaciones Específicas por Plataforma}
\label{tab:platform-optimization}
\small
\begin{tabular}{|l|p{9cm}|}
\hline
\textbf{Plataforma} & \textbf{Optimizaciones Requeridas} \\
\hline
Web Desktop & Lazy loading imágenes, code splitting por ruta, service workers para caché \\
\hline
Web Mobile & Reducción de bundle, imágenes optimizadas (WebP), evitar layout shifts \\
\hline
iOS (Futuro) & Reducir uso de memoria, caché de imágenes, optimizar animaciones Core Animation \\
\hline
Android (Futuro) & Soportar múltiples densidades (mdpi-xxxhdpi), reducir APK size, ProGuard \\
\hline
\end{tabular}
\end{table}

\section{Lineamientos de Contenido y Comunicación}

Esta sección establece los principios de redacción y tono comunicativo que deben emplearse en todas las interacciones textuales del sistema. La voz de marca coherente refuerza la identidad organizacional y construye relación con los usuarios (Redish, 2014).

\subsection{Tono y Voz de Marca}

El tono de comunicación de FINGEST debe balancear profesionalismo con accesibilidad, generando confianza sin intimidar. En el contexto de gestión financiera, los usuarios esperan seriedad y competencia, pero también claridad y soporte.

Las características del tono incluyen: profesional pero cercano evitando formalismo excesivo, claro y directo minimizando jerga financiera innecesaria, optimista enfocándose en logros y crecimiento, y empático reconociendo las preocupaciones financieras legítimas de los usuarios. El lenguaje debe utilizar voz activa preferentemente, segunda persona (tú/usted según contexto regional), frases concisas con máximo 20 palabras, y terminología consistente para conceptos recurrentes.

\subsection{Microcopy y Mensajes del Sistema}

El microcopy incluye todos los textos breves de la interfaz: botones, enlaces, placeholders, mensajes de error, y confirmaciones. Estos textos pequeños tienen impacto desproporcionado en la usabilidad y percepción de calidad (Kinneret, 2017).

Los botones deben utilizar verbos de acción específicos describiendo el resultado: ``Guardar cambios'', ``Crear cuenta'', ``Transferir dinero'' para acciones primarias; ``Cancelar'', ``Volver'', ``Omitir'' para acciones secundarias; y ``Eliminar transacción'', ``Cerrar cuenta'' para acciones destructivas. Los mensajes de error deben ser específicos explicando exactamente qué salió mal (``Ingresa un correo electrónico válido'' no ``Error en el campo''), constructivos indicando cómo corregir el problema, sin culpar al usuario evitando frases como ``Tu error fue'', y cuando sea posible, ofrecer solución automática mediante botón de acción.

Las notificaciones clasificadas por tipo utilizan lenguaje apropiado: éxito con ``¡Listo! Tu transacción se guardó correctamente'', información con ``Estamos procesando tu solicitud. Te notificaremos cuando esté lista'', advertencia con ``Tu sesión expirará en 5 minutos. ¿Deseas continuar?'', y error con ``No pudimos procesar el pago. Verifica los datos de tu tarjeta e intenta nuevamente''. Los placeholders en campos de entrada deben mostrar ejemplo del formato esperado (``nombre@ejemplo.com'') en lugar de instrucciones genéricas (``Ingresa tu correo'').

\section{Aplicación de Principios Teóricos en Casos Prácticos}

Esta sección mapea fundamentos teóricos (Gestalt, Nielsen, WCAG) a implementaciones específicas en FINGEST mediante tablas técnicas de referencia.

\subsection{Implementación de Principios de Gestalt}

Los principios perceptuales de Gestalt se aplican sistemáticamente. La Tabla \ref{tab:gestalt-implementacion} detalla cada principio con su ubicación específica.

\begin{table}[h]
\centering
\caption{Implementación de Principios Gestalt en FINGEST}
\label{tab:gestalt-implementacion}
\scriptsize
\begin{tabular}{|p{2cm}|p{2.5cm}|p{4cm}|p{3.5cm}|}
\hline
\textbf{Principio} & \textbf{Ubicación} & \textbf{Implementación Técnica} & \textbf{Código/Especificación} \\
\hline
\multirow{4}{*}{\textbf{Proximidad}} 
& Login: Campos relacionados 
& Email + Password agrupados, margin-bottom: 8px entre ellos, 24px antes del botón 
& \texttt{<div class="space-y-2">} \\
\cline{2-4}
& Dashboard: Métricas 
& Cards adyacentes, gap: 16px, agrupadas en grid 
& \texttt{grid grid-cols-3 gap-4} \\
\cline{2-4}
& Transacciones: Filtros 
& Filtros en bloque superior, padding: 16px, border-bottom separador 
& \texttt{<div class="filters-container">} \\
\cline{2-4}
& Formularios: Secciones 
& Grupos de campos, margin-bottom: 32px entre secciones 
& \texttt{<fieldset class="mb-8">} \\
\hline
\multirow{4}{*}{\textbf{Similitud}} 
& Sidebar: Enlaces 
& Todos los items: height: 44px, padding: 12px 16px 
& \texttt{.nav-item \{ height: 44px; \}} \\
\cline{2-4}
& Cards de información 
& Border: 1px \#E5E1DD, radius: 8px, background: \#FFF 
& \texttt{<Card class="border rounded-lg">} \\
\cline{2-4}
& Acciones primarias 
& Todos los botones principales: background: \#227C91 
& \texttt{<Button variant="primary">} \\
\cline{2-4}
& Iconos de categoría 
& Mismo tamaño (20px), stroke: 2px 
& \texttt{<Icon size=\{20\} />} \\
\hline
\multirow{4}{*}{\textbf{Continuidad}} 
& Formularios verticales 
& Campos alineados, width: 100\%, labels arriba 
& \texttt{<form class="space-y-4">} \\
\cline{2-4}
& Grid de métricas 
& Cards en línea horizontal, mismo altura 
& \texttt{grid grid-cols-4 gap-4} \\
\cline{2-4}
& Breadcrumbs 
& Separador chevron (›), items en línea 
& \texttt{<Breadcrumbs separator="›">} \\
\cline{2-4}
& Tablas de datos 
& Líneas horizontales entre filas 
& \texttt{tr \{ border-bottom: 1px solid; \}} \\
\hline
\end{tabular}
\end{table}

\newpage

\subsection{Aplicación de Heurísticas de Nielsen}

Las 10 heurísticas de Nielsen (1994) se implementan en elementos específicos. La Tabla \ref{tab:nielsen-implementacion} mapea cada heurística.


\begin{table}[h]
\centering
\caption{Implementación de Heurísticas de Nielsen en FINGEST}
\label{tab:nielsen-implementacion}
\scriptsize
\setlength{\tabcolsep}{4pt}
\begin{tabular}{|p{3cm}|p{2cm}|p{5.2cm}|}
\hline
\textbf{Heurística} & \textbf{Ubicación} & \textbf{Implementación} \\
\hline
\multirow{2}{*}{1. Estado del sistema} 
& Confirmación acciones & Toast notifications, duración 3 segundos, posición top-right \\
\cline{2-3}
& Carga asíncrona & Spinner 32px color primario, barra de progreso para operaciones largas \\
\hline
\multirow{2}{*}{2. Sistema-mundo real} 
& Términos financieros & Labels familiares: "Balance", "Transacciones", "Transferencias" \\
\cline{2-3}
& Iconografía & Iconos reconocibles: CreditCard, TrendingUp, Wallet, ShoppingCart \\
\hline
\multirow{2}{*}{3. Control y libertad} 
& Filtros & Botón "Limpiar filtros" para restaurar estado inicial \\
\cline{2-3}
& Operaciones & Modal con botón "Cancelar" siempre visible, navegación libre \\
\hline
\multirow{2}{*}{4. Estándares} 
& Componentes UI & Biblioteca consistente de botones, inputs, cards reutilizables \\
\cline{2-3}
& Interacciones & Patrones consistentes: hover, focus, active states uniformes \\
\hline
\multirow{2}{*}{5. Prevención errores} 
& Validación & Check en tiempo real: email pattern, contraseña mínima 8 chars \\
\cline{2-3}
& Confirmación & Modal "¿Estás seguro?" antes de eliminar, botones deshabilitados \\
\hline
\multirow{2}{*}{6. Reconocimiento} 
& Autocompletado & Dropdown de sugerencias, valores recientes primero \\
\cline{2-3}
& Valores default & Campos pre-llenados con datos comunes o del usuario \\
\hline
\multirow{2}{*}{7. Eficiencia} 
& Atajos teclado & Ctrl+S guardar, Esc cerrar modales, Ctrl+K buscar global \\
\cline{2-3}
& Búsqueda rápida & Comand palette accesible, filtros avanzados colapsables \\
\hline
\multirow{2}{*}{8. Minimalista} 
& Espaciado & Padding generoso (24-32px), espacios en blanco intencionales \\
\cline{2-3}
& Tipografía & Fuente Inter, pesos limitados, jerarquía clara \\
\hline
\multirow{2}{*}{9. Ayuda errores} 
& Mensajes & Específicos y descriptivos: "Email inválido: falta @" \\
\cline{2-3}
& Sugerencias & Texto de ayuda: "Intenta: usuario@example.com" \\
\hline
\multirow{2}{*}{10. Documentación} 
& Tooltips & Explicaciones breves al hover sobre iconos complejos \\
\cline{2-3}
& Helper text & Instrucciones debajo de campos: formato esperado \\
\hline
\end{tabular}
\end{table}

\newpage

\subsection{Matriz de Validación de Principios}

\begin{table}[h]
\centering
\caption{Matriz de Cobertura de Principios por Componente}
\label{tab:matriz-validacion}
\scriptsize
\begin{tabular}{|l|c|c|c|c|c|c|}
\hline
\textbf{Componente} & \textbf{Gestalt} & \textbf{Nielsen} & \textbf{WCAG A} & \textbf{WCAG AA} & \textbf{WCAG AAA} & \textbf{Cobertura} \\
\hline
Botones & ✓ & H1, H4, H5 & ✓ & ✓ & ✓ & 100\% \\
\hline
Inputs & ✓ & H1, H5, H9 & ✓ & ✓ & — & 95\% \\
\hline
Cards & ✓ & H4, H8 & ✓ & ✓ & — & 90\% \\
\hline
Navegación & ✓ & H2, H3, H4 & ✓ & ✓ & ✓ & 100\% \\
\hline
Formularios & ✓ & H5, H6, H9 & ✓ & ✓ & — & 100\% \\
\hline
Modales & ✓ & H1, H3, H5 & ✓ & ✓ & — & 95\% \\
\hline
Tablas & ✓ & H4, H8 & ✓ & ✓ & — & 90\% \\
\hline
Notificaciones & ✓ & H1, H9 & ✓ & ✓ & — & 100\% \\
\hline
\end{tabular}
\end{table}

\textbf{Leyenda:} H1-H10 = Heurísticas de Nielsen (1-10), ✓ = Implementado, — = No aplica


\section{Gestión y Evolución del Sistema de Diseño}

Esta sección establece los procesos organizacionales necesarios para mantener la calidad y relevancia del manual de identidad visual a lo largo del tiempo. Un sistema de diseño efectivo no es estático, sino que evoluciona respondiendo a nuevas necesidades, tecnologías y feedback de usuarios.

\subsection{Proceso de Revisión y Actualización}

La revisión sistemática del manual asegura que las directrices permanezcan relevantes y efectivas. El proceso involucra múltiples niveles de responsabilidad donde el diseñador principal actúa como custodio del sistema, el equipo de desarrollo proporciona feedback sobre viabilidad técnica, el área de QA valida cumplimiento de estándares, y el product owner aprueba cambios significativos que afecten la dirección estratégica.

La frecuencia de revisión se estructura en ciclos trimestrales para evaluación general de adherencia y consistencia, revisiones post-release después de implementar funcionalidades principales que puedan haber revelado necesidades nuevas, y revisiones ad-hoc cuando se identifiquen inconsistencias críticas o problemas de accesibilidad. Cada revisión debe documentarse con fecha, participantes, hallazgos principales, y acciones recomendadas.

El proceso de cambios formales incluye: (a) propuesta documentada con justificación basada en datos de usuarios o limitaciones técnicas, (b) revisión colaborativa por equipo de diseño evaluando impacto en coherencia visual, (c) creación de prototipos o mockups del cambio propuesto, (d) testing de accesibilidad y usabilidad con usuarios reales cuando sea posible, (e) aprobación formal documentada, y (f) implementación gradual con monitoreo de impacto.

\subsection{Sistema de Versionado}

El versionado semántico (Major.Minor.Patch) proporciona claridad sobre la magnitud de cambios realizados. Las versiones Major (1.0.0 → 2.0.0) indican cambios fundamentales de identidad visual como nueva paleta de colores, rediseño completo de componentes principales, o cambio de tipografía corporativa. Estos cambios requieren aprobación ejecutiva y planificación de migración.

Las versiones Minor (1.0.0 → 1.1.0) introducen nuevos componentes o guidelines sin romper la compatibilidad existente, como adición de nuevos patrones de interacción, expansión de paleta de colores con tonos complementarios, o documentación de casos de uso nuevos. Las versiones Patch (1.1.0 → 1.1.1) corrigen errores, aclaran ambigüedades, o actualizan ejemplos sin modificar especificaciones fundamentales.

Cada versión debe documentarse en archivo CHANGELOG.md incluyendo fecha de publicación, número de versión, resumen ejecutivo de cambios, lista detallada de adiciones, modificaciones y correcciones, referencias a issues o pull requests relacionados, screenshots comparativos cuando sea relevante, y notas de migración para cambios que requieran actualización de código existente.

\section{Conclusiones}

Este manual de identidad visual y experiencia de usuario establece un marco integral para el desarrollo coherente de productos digitales bajo la marca FINGEST. La sistematización de elementos visuales, patrones de interacción y principios de diseño presentados en este documento no constituye simplemente un catálogo de especificaciones técnicas, sino una filosofía de diseño centrada en el usuario que equilibra estética, funcionalidad y accesibilidad.

La integración de fundamentos teóricos sólidos (Gestalt, Nielsen, WCAG) con implementaciones prácticas asegura que las decisiones de diseño no sean arbitrarias sino fundamentadas en investigación científica y mejores prácticas de la industria. La paleta cromática cumple rigurosamente con estándares de accesibilidad WCAG 2.1 nivel AA y AAA, el sistema tipográfico optimiza legibilidad en múltiples dispositivos y contextos, los componentes UI reutilizables aceleran el desarrollo mientras mantienen consistencia, y los patrones de interacción documentados facilitan experiencias predecibles y eficientes.

La implementación exitosa de estas directrices requiere compromiso organizacional más allá del equipo de diseño. Desarrolladores, gerentes de producto, profesionales de QA y marketing deben comprender y aplicar estos principios en su trabajo diario. La consistencia visual genera confianza, elemento fundamental en aplicaciones de gestión financiera donde los usuarios depositan información sensible y toman decisiones económicas importantes.

Las recomendaciones para fortalecimiento continuo incluyen: auditoría sistemática de cumplimiento en componentes existentes, desarrollo de biblioteca Storybook para documentación interactiva accesible, completar activos visuales faltantes (logotipo SVG profesional, sistema iconográfico expandido), implementar testing automatizado de accesibilidad en pipelines de CI/CD, establecer proceso de revisión trimestral con participación multi-disciplinaria, planificar expansión a plataformas móviles nativas respetando convenciones de cada ecosistema, y fomentar cultura de feedback continuo donde todos los miembros del equipo se sientan empoderados para sugerir mejoras.

El sistema de diseño es un documento vivo que debe evolucionar respondiendo a necesidades cambiantes de usuarios, avances tecnológicos, y aprendizajes derivados de la operación real del sistema. Se invita a todos los equipos a contribuir activamente reportando inconsistencias, proponiendo componentes nuevos, documentando casos de uso emergentes, y participando en sesiones de capacitación y revisión periódicas. Solo mediante iteración continua y mejora basada en evidencia el manual mantendrá su relevancia y utilidad a largo plazo.


\section{Referencias}

\begin{flushleft}

Bernard, M., Liao, C. H., \& Mills, M. (2003). The effects of font type and size on the legibility and reading time of online text by older adults. \textit{CHI'03 Extended Abstracts on Human Factors in Computing Systems}, 175-176.

\vspace{0.5cm}

Galitz, W. O. (2007). \textit{The Essential Guide to User Interface Design: An Introduction to GUI Design Principles and Techniques} (3rd ed.). Wiley Publishing.

\vspace{0.5cm}

Henry, S. L., Abou-Zahra, S., \& Brewer, J. (2014). The role of accessibility in a universal web. \textit{Proceedings of the 11th Web for All Conference}, Article 17, 1-4.

\vspace{0.5cm}

Huang, S. M., Shieh, K. K., \& Chi, C. F. (2002). Factors affecting the design of computer icons. \textit{International Journal of Industrial Ergonomics}, 29(4), 211-218.

\vspace{0.5cm}

Jarrett, C., \& Gaffney, G. (2009). \textit{Forms that Work: Designing Web Forms for Usability}. Morgan Kaufmann Publishers.

\vspace{0.5cm}

Kalbach, J. (2007). \textit{Designing Web Navigation: Optimizing the User Experience}. O'Reilly Media.

\vspace{0.5cm}

Kinneret, Y. (2017). \textit{Microcopy: The Complete Guide}. Apress.

\vspace{0.5cm}

Koffka, K. (1935). \textit{Principles of Gestalt Psychology}. Harcourt, Brace \& World.

\vspace{0.5cm}

Krug, S. (2014). \textit{Don't Make Me Think, Revisited: A Common Sense Approach to Web Usability} (3rd ed.). New Riders.

\vspace{0.5cm}

Labrecque, L. I., \& Milne, G. R. (2012). Exciting red and competent blue: The importance of color in marketing. \textit{Journal of the Academy of Marketing Science}, 40(5), 711-727.

\vspace{0.5cm}

Laubheimer, P. (2020). \textit{The Kinetics of Web Animations}. Nielsen Norman Group. Recuperado de https://www.nngroup.com/articles/animation-duration/

\vspace{0.5cm}

Marcotte, E. (2011). \textit{Responsive Web Design}. A Book Apart.

\vspace{0.5cm}

Nielsen, J. (1994). Heuristic evaluation. En J. Nielsen \& R. L. Mack (Eds.), \textit{Usability Inspection Methods} (pp. 25-62). John Wiley \& Sons.

\vspace{0.5cm}

Norman, D. A. (2013). \textit{The Design of Everyday Things: Revised and Expanded Edition}. Basic Books.

\vspace{0.5cm}

Palmer, S. E. (1992). Common region: A new principle of perceptual grouping. \textit{Cognitive Psychology}, 24(3), 436-447.

\vspace{0.5cm}

Redish, J. (2014). \textit{Letting Go of the Words: Writing Web Content that Works} (2nd ed.). Morgan Kaufmann.

\vspace{0.5cm}

Shaikh, A. D. (2007). Psychology of onscreen type: Investigations regarding typeface personality, appropriateness, and impact on document perception [Tesis doctoral]. Wichita State University.

\vspace{0.5cm}

Todorovic, D. (2008). Gestalt principles. \textit{Scholarpedia}, 3(12), 5345.

\vspace{0.5cm}

W3C. (2018). \textit{Web Content Accessibility Guidelines (WCAG) 2.1}. World Wide Web Consortium. Recuperado de https://www.w3.org/TR/WCAG21/

\vspace{0.5cm}

Wertheimer, M. (1923). Untersuchungen zur Lehre von der Gestalt II. \textit{Psychologische Forschung}, 4(1), 301-350.

\vspace{0.5cm}

Wheeler, A. (2017). \textit{Designing Brand Identity: An Essential Guide for the Whole Branding Team} (5th ed.). John Wiley \& Sons.

\end{flushleft}

\newpage


\section{Anexos}

\subsection{Anexo A: Especificaciones Técnicas de Colores}

Este anexo presenta las especificaciones completas de la paleta cromática de FINGEST, incluyendo códigos en múltiples formatos para facilitar implementación en diversos contextos técnicos. Cada color ha sido validado para cumplir con ratios de contraste WCAG 2.1.

\begin{center}
\begin{tabular}{|l|c|l|l|l|p{4cm}|}
\hline
\textbf{Nombre} & \textbf{Muestra} & \textbf{HEX} & \textbf{RGB} & \textbf{HSL} & \textbf{Uso Principal} \\
\hline
Primary & \colorbox{primary}{\color{white}\rule{0pt}{12pt}\rule{20pt}{0pt}} & \#227C91 & 34, 124, 145 & 191° 62\% 35\% & Botones CTA, enlaces principales \\
\hline
Secondary & \colorbox{secondary}{\color{white}\rule{0pt}{12pt}\rule{20pt}{0pt}} & \#605952 & 96, 89, 82 & 30° 8\% 35\% & Texto secundario, elementos neutros \\
\hline
Destructive & \colorbox{destructive}{\color{white}\rule{0pt}{12pt}\rule{20pt}{0pt}} & \#EF4444 & 239, 68, 68 & 0° 84\% 60\% & Errores, alertas, borrado \\
\hline
Background & \colorbox{background}{\color{black}\rule{0pt}{12pt}\rule{20pt}{0pt}} & \#FAFAFA & 250, 250, 250 & 40° 6\% 98\% & Fondo principal aplicación \\
\hline
Muted & \colorbox{muted}{\color{black}\rule{0pt}{12pt}\rule{20pt}{0pt}} & \#F1F0EE & 241, 240, 238 & 40° 7\% 94\% & Fondos deshabilitados \\
\hline
Border & \colorbox{border}{\color{black}\rule{0pt}{12pt}\rule{20pt}{0pt}} & \#E5E1DD & 229, 225, 221 & 30° 16\% 88\% & Bordes, separadores \\
\hline
\end{tabular}
\end{center}

\vspace{0.5cm}

\textbf{Ratios de Contraste (sobre fondo blanco):}
\begin{itemize}
    \item Primary: 7.8:1 (Nivel AAA)
    \item Secondary: 7.2:1 (Nivel AAA)
    \item Destructive: 4.7:1 (Nivel AA)
    \item Muted Foreground sobre Muted: 4.8:1 (Nivel AA)
\end{itemize}

\subsection{Anexo B: Jerarquía Tipográfica Completa}

Este anexo detalla el sistema tipográfico completo de FINGEST basado en la fuente Inter. Las especificaciones aseguran legibilidad óptima en dispositivos digitales y establecen jerarquía visual clara.

\begin{center}
\begin{tabular}{|l|l|l|l|p{5cm}|}
\hline
\textbf{Elemento} & \textbf{Tamaño} & \textbf{Peso} & \textbf{Line Height} & \textbf{Contexto de Uso} \\
\hline
H1 & 32px (2rem) & 700 (Bold) & 1.2 & Títulos de página principales, encabezados hero \\
\hline
H2 & 24px (1.5rem) & 700 (Bold) & 1.3 & Títulos de sección, divisiones mayores \\
\hline
H3 & 20px (1.25rem) & 600 (Semibold) & 1.4 & Subsecciones, títulos de tarjetas \\
\hline
Párrafo & 16px (1rem) & 400 (Regular) & 1.5 & Texto de cuerpo estándar, descripciones \\
\hline
Pequeño & 14px (0.875rem) & 400 (Regular) & 1.4 & Texto complementario, metadatos \\
\hline
Ayuda & 12px (0.75rem) & 400 (Regular) & 1.3 & Anotaciones, tooltips, notas al pie \\
\hline
\end{tabular}
\end{center}

\vspace{0.5cm}

\textbf{Consideraciones de Implementación:}
\begin{itemize}
    \item Cargar fuente desde Google Fonts: \texttt{@import url('https://fonts.googleapis.com/css2?family=Inter:wght@400;500;600;700\&display=swap')}
    \item Aplicar \texttt{font-family: 'Inter', sans-serif} en elemento raíz
    \item Letter-spacing: usar valores por defecto (Inter ya está optimizada)
    \item Utilizar \texttt{font-feature-settings} para activar ligaduras opcionales
\end{itemize}

\subsection{Anexo C: Checklist de Accesibilidad}

Este checklist proporciona una guía práctica para validar el cumplimiento de estándares de accesibilidad WCAG 2.1 en componentes de interfaz de FINGEST. Debe utilizarse en revisiones de código, testing de QA, y auditorías periódicas del sistema.

\textbf{Contraste de Color:}
\begin{itemize}
    \item[$\square$] Texto normal (menor a 18px) cumple ratio mínimo 4.5:1 (Nivel AA)
    \item[$\square$] Texto grande (18px+ o 14px+ en negrita) cumple ratio 3:1 (Nivel AA)
    \item[$\square$] Componentes de interfaz (botones, campos) cumplen ratio 3:1
    \item[$\square$] Información crítica no depende únicamente del color
    \item[$\square$] Estados (error, éxito, advertencia) usan iconos además de color
\end{itemize}

\textbf{Navegación por Teclado:}
\begin{itemize}
    \item[$\square$] Todos los elementos interactivos alcanzables con tecla Tab
    \item[$\square$] Orden de tabulación sigue flujo lógico (izquierda-derecha, arriba-abajo)
    \item[$\square$] Indicadores de focus visibles con contraste adecuado (anillo 2px)
    \item[$\square$] No existen trampas de teclado (elementos de los que no se puede salir)
    \item[$\square$] Tecla Escape cierra modales y overlays
    \item[$\square$] Enter/Space activan botones y elementos interactivos
\end{itemize}

\textbf{Contenido Semántico:}
\begin{itemize}
    \item[$\square$] Estructura HTML semántica apropiada (header, nav, main, footer, article)
    \item[$\square$] Headings en orden jerárquico sin saltos (H1 → H2 → H3, no H1 → H3)
    \item[$\square$] Imágenes funcionales tienen texto alternativo descriptivo
    \item[$\square$] Imágenes decorativas marcadas con \texttt{alt=""} o \texttt{aria-hidden="true"}
    \item[$\square$] Formularios usan \texttt{<label>} asociados correctamente con \texttt{htmlFor}
    \item[$\square$] Campos requeridos indicados claramente (asterisco + aria-required)
\end{itemize}

\textbf{Compatibilidad con Tecnologías Asistivas:}
\begin{itemize}
    \item[$\square$] Probado con NVDA (Windows) en Chrome y Firefox
    \item[$\square$] Probado con VoiceOver (Mac/iOS) en Safari
    \item[$\square$] Atributos ARIA utilizados apropiadamente donde HTML semántico es insuficiente
    \item[$\square$] Contenido dinámico usa \texttt{aria-live} regions para notificar cambios
    \item[$\square$] Componentes personalizados incluyen roles ARIA apropiados
    \item[$\square$] Estados de componentes comunicados con \texttt{aria-expanded}, \texttt{aria-selected}
\end{itemize}

\textbf{Multimedia y Animaciones:}
\begin{itemize}
    \item[$\square$] Animaciones respetan preferencia \texttt{prefers-reduced-motion}
    \item[$\square$] Videos incluyen subtítulos y/o transcripciones
    \item[$\square$] Audio no reproduce automáticamente o puede pausarse fácilmente
    \item[$\square$] Animaciones pueden pausarse, detenerse u ocultarse
\end{itemize}

\subsection{Anexo D: Breakpoints y Adaptación Responsiva}

Este anexo especifica los breakpoints y estrategias de adaptación para diferentes tamaños de pantalla. El diseño responsivo de FINGEST sigue el enfoque mobile-first, construyendo progresivamente desde dispositivos pequeños hacia pantallas grandes.

\begin{center}
\begin{tabular}{|l|l|p{3cm}|p{6cm}|}
\hline
\textbf{Breakpoint} & \textbf{Rango} & \textbf{Dispositivos Típicos} & \textbf{Adaptaciones Principales} \\
\hline
xs (default) & 0-639px & Smartphones verticales & Sidebar oculto (drawer), grid 1 columna, navegación simplificada \\
\hline
sm & 640-767px & Smartphones horizontales & Grids selectivos 1-2 columnas, tipografía ligeramente mayor \\
\hline
md & 768-1023px & Tablets verticales & Sidebar colapsable a iconos, grids 2 columnas, formularios expandidos \\
\hline
lg & 1024-1279px & Tablets horizontales, laptops & Sidebar visible completo, grids 3 columnas, layout multi-panel \\
\hline
xl & 1280px+ & Monitores desktop & Layout completo, uso óptimo espacio horizontal, paneles laterales \\
\hline
\end{tabular}
\end{center}

\vspace{0.5cm}

\textbf{Estrategias de Adaptación por Componente:}

\textit{Navegación:}
\begin{itemize}
    \item xs-sm: Menú hamburguesa, drawer deslizable desde izquierda
    \item md: Sidebar colapsado mostrando solo iconos (opción de expandir)
    \item lg-xl: Sidebar completo visible permanentemente (256px ancho)
\end{itemize}

\textit{Grids y Layouts:}
\begin{itemize}
    \item xs: Una columna apilada verticalmente
    \item sm: Contenido flexible, hasta 2 columnas en secciones anchas
    \item md: Dos columnas estándar, tres columnas en secciones específicas
    \item lg-xl: Tres columnas estándar, layouts complejos multi-panel
\end{itemize}

\textit{Formularios:}
\begin{itemize}
    \item xs-sm: Campos apilados en una columna
    \item md-xl: Campos en dos columnas donde sea lógico (nombre/apellido)
    \item Botones de acción siempre accesibles (sticky en móvil si es necesario)
\end{itemize}

\textit{Tablas de Datos:}
\begin{itemize}
    \item xs-sm: Transformar a vistas de tarjetas apiladas o scroll horizontal
    \item md: Tabla estándar con columnas prioritarias visibles
    \item lg-xl: Tabla completa con todas las columnas
\end{itemize}

\textbf{Testing en Dispositivos Reales:}

Se recomienda probar regularmente en:
\begin{itemize}
    \item iPhone SE (320px) - Límite inferior de soporte
    \item iPhone 13/14 (390px) - Smartphone estándar actual
    \item iPad Mini (768px) - Tablet pequeña vertical
    \item iPad Pro (1024px) - Tablet grande horizontal
    \item MacBook Pro 13" (1440px) - Laptop estándar
    \item Monitor 4K (2560px+) - Desktop grande
\end{itemize}

\newpage

\subsection{Anexo E: Principios de Diseño aplicados en FINGEST}

% ------------------------------------------------------------
\subsubsection{E.1 Interfaz de Registro de Usuario}

\begin{figure}[h]
    \centering
    \includegraphics[width=0.65\linewidth]{register.png}
    \caption{Interfaz de Registro de Usuario en FINGEST}
    \label{fig:register}
\end{figure}

\noindent \textbf{Principios aplicados:}

\subsectionautorefname{E.1.1 Principios de Nielsen (Heurísticas de Usabilidad)}
\begin{itemize}
    \item \textbf{Visibilidad del estado del sistema:}  
    El usuario recibe retroalimentación inmediata al ingresar datos o cometer errores (por ejemplo, campos requeridos o contraseñas no válidas).

    \item \textbf{Correspondencia entre el sistema y el mundo real:}  
    Se usan etiquetas claras como "Correo electrónico" y "Contraseña", evitando tecnicismos.

    \item \textbf{Control y libertad del usuario:}  
    Los botones "Cancelar" o "Regresar" permiten salir sin perder el control de la acción.

    \item \textbf{Consistencia y estándares:}  
    El diseño de formularios sigue el mismo estilo que el resto del sistema: campos con bordes redondeados, tipografía clara y colores consistentes.

    \item \textbf{Prevención de errores:}  
    Se validan los datos antes del envío, mostrando advertencias visuales cuando algo no está completo.

    \item \textbf{Estética y diseño minimalista:}  
    La pantalla de registro es simple, priorizando la claridad visual y evitando distracciones.
\end{itemize}

\subsubsection*{E.1.2 Principios de la Gestalt}
\begin{itemize}
    \item \textbf{Ley de proximidad:}  
    Los campos de texto y botones están agrupados según su función (datos personales, botones de acción).

    \item \textbf{Ley de similitud:}  
    Todos los campos tienen el mismo formato visual, facilitando la identificación del patrón de llenado.

    \item \textbf{Ley de figura-fondo:}  
    El formulario resalta sobre un fondo neutro, guiando la atención del usuario al área principal.
\end{itemize}

% ------------------------------------------------------------
\subsubsection{E.2 Panel Principal o Dashboard del Sistema}

\begin{figure}[h]
    \centering
    \includegraphics[width=0.65\linewidth]{dashboard.png}
    \caption{Panel Principal o Dashboard del Sistema}
    \label{fig:dashboard}
\end{figure}

\noindent \textbf{Principios aplicados:}

\subsubsection*{E.2.1 Principios de Nielsen (Heurísticas de Usabilidad)}
\begin{itemize}
    \item \textbf{Visibilidad del estado del sistema:}  
    Muestra de forma inmediata el total de ingresos, gastos y balance neto.

    \item \textbf{Correspondencia entre el sistema y el mundo real:}  
    Se usan términos familiares (por ejemplo, "Transacciones", "Gastos", "Ingresos").

    \item \textbf{Consistencia y estándares:}  
    Cada tarjeta utiliza los mismos colores, iconos y disposición, manteniendo coherencia visual.

    \item \textbf{Reconocimiento mejor que recuerdo:}  
    Los íconos y colores indican el tipo de operación (verde para ingresos, rojo para gastos).

    \item \textbf{Estética y diseño minimalista:}  
    Diseño limpio, sin información innecesaria, con foco en los datos financieros principales.
\end{itemize}

\subsubsection*{E.2.2 Principios de la Gestalt}
\begin{itemize}
    \item \textbf{Ley de proximidad:}  
    Los valores de ingresos, gastos y balance se agrupan visualmente para facilitar la comparación.

    \item \textbf{Ley de similitud:}  
    Las tarjetas comparten estructura visual similar, lo que genera orden y facilita la lectura.

    \item \textbf{Ley de figura-fondo:}  
    Las tarjetas blancas destacan sobre el fondo gris, facilitando la concentración del usuario.
\end{itemize}

 \newpage

% ------------------------------------------------------------
\subsubsection{E.3 Barra Lateral de Navegación del Sistema}

\begin{figure}[h]
    \centering
    \includegraphics[height=0.40\textheight, keepaspectratio]{sidebar.png}
    \caption{Barra Lateral de Navegación del Sistema}
    \label{fig:sidebar}
\end{figure}

\noindent \textbf{Principios aplicados:}

\subsubsection*{E.3.1 Principios de Nielsen (Heurísticas de Usabilidad)}
\begin{itemize}
    \item \textbf{Consistencia y estándares:}  
    Los iconos y etiquetas mantienen un estilo uniforme y reconocible.

    \item \textbf{Visibilidad del estado del sistema:}  
    El elemento activo se resalta para indicar en qué sección se encuentra el usuario.

    \item \textbf{Correspondencia con el mundo real:}  
    Iconos como "��", "��" o "��" representan acciones familiares (inicio, trabajo, reportes).

    \item \textbf{Flexibilidad y eficiencia de uso:}  
    La estructura fija de la barra lateral permite navegación rápida sin pérdida de contexto.

    \item \textbf{Estética y diseño minimalista:}  
    Presenta solo los elementos esenciales, con iconografía limpia.
\end{itemize}

\subsubsection*{E.3.2 Principios de la Gestalt}
\begin{itemize}
    \item \textbf{Ley de continuidad:}  
    La disposición vertical guía el recorrido visual de forma natural.

    \item \textbf{Ley de similitud:}  
    Todos los botones mantienen la misma forma y color, reforzando la coherencia visual.
\end{itemize}

% ------------------------------------------------------------
\subsubsection{E.4 Interfaz de Eventos y Actividades Financieras}

\begin{figure}[h]
    \centering
    \includegraphics[width=0.65\linewidth]{noti_eventos.png}
    \caption{Interfaz de Eventos y Actividades Financieras}
    \label{fig:noti_eventos}
\end{figure}

\noindent \textbf{Principios aplicados:}

\subsubsection*{E.4.1 Principios de Nielsen (Heurísticas de Usabilidad)}
\begin{itemize}
    \item \textbf{Visibilidad del estado del sistema:}  
    Las tarjetas muestran estados como "Evento activo", "Próximo pago" o "Finalizado".

    \item \textbf{Correspondencia entre el sistema y el mundo real:}  
    Se usan fechas y etiquetas realistas (por ejemplo, "Jul 20, 2023 – Taller de presupuesto").

    \item \textbf{Consistencia y estándares:}  
    Cada evento sigue un formato idéntico de información: ícono, título, fecha y acción.

    \item \textbf{Prevención de errores:}  
    Las alertas anticipan plazos o pagos próximos, evitando olvidos.

    \item \textbf{Estética y diseño minimalista:}  
    Colores suaves y separación clara entre eventos mejoran la legibilidad.
\end{itemize}

\subsubsection*{E.4.2 Principios de la Gestalt}
\begin{itemize}
    \item \textbf{Ley de proximidad:}  
    Los elementos relacionados dentro de cada evento se agrupan en una sola tarjeta.

    \item \textbf{Ley de continuidad:}  
    El flujo visual permite recorrer la lista de eventos de forma ordenada y fluida.
\end{itemize}

% ------------------------------------------------------------
\subsubsection{E.5 Vista de Notificaciones del Sistema Financiero}

\begin{figure}[h]
    \centering
    \includegraphics[width=0.75\linewidth]{notoficaciones.png}
    \caption{Vista de Notificaciones del Sistema Financiero}
    \label{fig:notificaciones}
\end{figure}

\noindent \textbf{Principios aplicados:}

\subsubsection*{E.5.1 Principios de Nielsen (Heurísticas de Usabilidad)}
\begin{itemize}
    \item \textbf{Visibilidad del estado del sistema:}  
    Cada tarjeta informa el estado actual ("Balance bajo", "Meta alcanzada", etc.) junto con fechas y acciones.

    \item \textbf{Correspondencia entre el sistema y el mundo real:}  
    Se usa lenguaje cotidiano y íconos familiares (⚠️, ✅, ��) para representar conceptos financieros reales.

    \item \textbf{Control y libertad del usuario:}  
    Los botones ❌ y "Marcar como leída" ofrecen control sobre las notificaciones.

    \item \textbf{Consistencia y estándares:}  
    Todas las tarjetas mantienen formato, colores y tipografía uniformes.

    \item \textbf{Prevención de errores:}  
    Se incluyen alertas como "Balance bajo" o "Gasto inusual" que ayudan a evitar errores.

    \item \textbf{Reconocimiento mejor que recuerdo:}  
    Los íconos y colores permiten identificar categorías sin leer todo el texto.

    \item \textbf{Flexibilidad y eficiencia de uso:}  
    Las pestañas ("Todas", "Importantes", "Recordatorios", "Metas", "Gastos") facilitan la personalización.

    \item \textbf{Estética y diseño minimalista:}  
    Interfaz clara y limpia, con jerarquía visual definida.
\end{itemize}

\subsubsection*{E.5.2 Principios de la Gestalt}
\begin{itemize}
    \item \textbf{Ley de proximidad:}  
    Íconos, texto y botones se agrupan en cada tarjeta formando unidades coherentes.

    \item \textbf{Ley de similitud:}  
    Estructura y estilo visual homogéneos en todas las tarjetas.

    \item \textbf{Ley de continuidad:}  
    Flujo vertical que facilita la lectura secuencial.

    \item \textbf{Ley de figura-fondo:}  
    Tarjetas resaltan sobre el fondo gris claro, diferenciando tipos de mensaje.

    \item \textbf{Ley de pregnancia (buena forma):}  
    Formas simples e íconos claros generan una percepción ordenada y agradable.
\end{itemize}

\end{document}
